\section{Further homotopy and detection results}
\section{Classical classes and synthetic comparison formulas}
\begin{definition}[Algebras on finite lambda quotients]
\label{frontier:Literature:63eb956dffc1}
\level{fine}
(Adapted from \cite{LWX126}, Remark 3.4.)\par{}

For each positive integer $q$, choose an algebra structure on $S^{0,0}/\lambda^q$ in the synthetic monoidal category.
\lean{KIP126.Literature.Standard.quotientAlgebra}
% Blueprint note: Source verification (2026-10-18): Source adaptation checked. AimPaper Remark 3.4 explicitly covers all positive powers. BHSmot Lemma B.5 directly covers the first special fiber and C.22 a completed cellular realization; it is not a standalone all-q theorem with the project’s exact choices. Primary references: [LWX126](https://arxiv.org/abs/2412.10879): Remark 3.4; pages 15; Repository-frozen AimPaper/main.tex and AimPaper/2412.10879.pdf [BHSmot](https://arxiv.org/pdf/2010.10325v2): Definition B.3; Lemma B.5; Proposition C.2; Proposition C.22; pages 69–70, 75; arXiv:2010.10325v2, 30 June 2022 [Pst](https://arxiv.org/pdf/1803.01804v3): Corollary 4.30; Theorem 4.46; Lemma 4.56; pages 46, 50–51, 54; arXiv:1803.01804v3, 10 November 2022
% Blueprint note: Detailed source attribution: (Adapted from \cite{LWX126}, Remark 3.4; see \cite{BHSmot}, Lemma B.5 and Proposition C.22.)
\end{definition}

\begin{definition}[The symmetric synthetic tensor structure]
\label{frontier:Literature:73db85fe3f01}
\level{fine}
(Adapted from \cite{Pst}, Prop. 4.2.)\par{}

Equip the synthetic category with its symmetric monoidal structure.
\lean{KIP126.Literature.Standard.quotientSyntheticSymmetric}
% Blueprint note: Source verification (2026-10-18): Source adaptation checked. The precise categorical source is Pstragowski Proposition 4.2. Recognition of the project’s selected synthetic model and its already fixed tensor remains pending. Primary references: [Pst](https://arxiv.org/pdf/1803.01804v3): Proposition 4.2; Remark 4.5; Definition 4.6; pages 38–39; arXiv:1803.01804v3, 10 November 2022
% Blueprint note: Detailed source attribution: (Adapted from \cite{Pst}, Proposition 4.2.)
\end{definition}

\begin{definition}[The classical bar-kappa class]
\label{frontier:Literature:051dbca71e3a}
\level{fine}
(Adapted from \cite{BR21}, Thm. 11.61(20).)\par{}

Fix the class $\bar\kappa\in\pi_{20}S^0$.
\lean{KIP126.Literature.Standard.kappaBar}
% Blueprint note: Source verification (2026-10-18): Source adaptation checked. The cited papers concern the usual classical spectrum or its 2-primary/2-completed version. Recognition of the fixed project category, Adams tower, coefficient object, and selected Lean constants is a separate, currently pending formal comparison; this audit does not prove it. Primary references: [BR21](https://www.sas.rochester.edu/mth/sites/doug-ravenel/otherpapers/rognes-bruner-tmf-book.pdf): Theorem11.61,case(20),p467;Table11.3,p458; Author preliminary version of AMS Mathematical Surveys and Monographs 253 (2021), downloaded 2026-10-18; printed page numbers cited, PDF index = printed page + 17 in the body.
% Blueprint note: Detailed source attribution: (Adapted from \cite{BR21}, Theorem 11.61(20), p. 467, and Table 11.3.)
\end{definition}



\begin{definition}[The Ext class g]
\label{frontier:Literature:174e2bfc0dfb}
\level{fine}
(Adapted from \cite{BR21}, Prop. 1.4.)\par{}

Fix $g\in E_2^{4,24}(S^0)$.
\lean{KIP126.Literature.Standard.g}
% Blueprint note: Source verification (2026-10-18): Source adaptation checked. The cited papers concern the usual classical spectrum or its 2-primary/2-completed version. Recognition of the fixed project category, Adams tower, coefficient object, and selected Lean constants is a separate, currently pending formal comparison; this audit does not prove it. Primary references: [BR21](https://www.sas.rochester.edu/mth/sites/doug-ravenel/otherpapers/rognes-bruner-tmf-book.pdf): Proposition1.4;Table1.1,p54, row(stem20,filtration4), cocycle4_8;Figure1.9,p56; Author preliminary version of AMS Mathematical Surveys and Monographs 253 (2021), downloaded 2026-10-18; printed page numbers cited, PDF index = printed page + 17 in the body.
% Blueprint note: Detailed source attribution: (Adapted from \cite{BR21}, Proposition 1.4 and Table 1.1, p. 54, in bidegree $(s,t)=(4,24)$.)
\end{definition}

\begin{definition}[The tmf source class]
\label{frontier:Literature:19f82a11f795}
\level{fine}
(Adapted from \cite{BR21}, Table 5.4.)\par{}

Fix the class $v_2^{16}\in E_2^{16,112}(\mathit{tmf})$.
\lean{KIP126.Literature.Standard.tmfV2Sixteen}
% Blueprint note: Source verification (2026-10-18): Source adaptation checked. The cited papers concern the usual classical spectrum or its 2-primary/2-completed version. Recognition of the fixed project category, Adams tower, coefficient object, and selected Lean constants is a separate, currently pending formal comparison; this audit does not prove it. Primary references: [BR21](https://www.sas.rochester.edu/mth/sites/doug-ravenel/otherpapers/rognes-bruner-tmf-book.pdf): BR21 Proposition1.12;Table1.3,p65,restrictionw2↦v2^8;Table5.4,p196,roww2²; AimPaper Appendix,paragraph immediately beforeTable5: d3(v2^16)=β^5g; Author preliminary version of AMS Mathematical Surveys and Monographs 253 (2021), downloaded 2026-10-18; printed page numbers cited, PDF index = printed page + 17 in the body. [LWX126](https://arxiv.org/abs/2412.10879): BR21 Proposition1.12;Table1.3,p65,restrictionw2↦v2^8;Table5.4,p196,roww2²; AimPaper Appendix,paragraph immediately beforeTable5: d3(v2^16)=β^5g; Repository AimPaper/main.tex and2412.10879.pdf, frozen baseline. Used only as secondary notation/recall; original BR21 and BMQ checked independently.
% Blueprint note: Detailed source attribution: (Adapted from \cite{BR21}, Proposition 1.12, Table 1.3, and Table 5.4: $v_2^{16}=w_2^2$ in the notation of \cite{LWX126}, Appendix.)
\end{definition}

\begin{definition}[The cofiber of a normalized arrow]
\label{frontier:Literature:1aa6a10db546}
\level{fine}
(Adapted from \cite{LWX126}, Prop. 3.20.)\par{}

For each classical triangle $X\xrightarrow{f}Y\xrightarrow{g}Z\to\Sigma X$ in the chosen family, let $e(g)$ denote the normalized filtration exponent. There is an isomorphism\[\operatorname{cofib}(\widehat f)\cong\Sigma^{0,-e(g)}\nu Z,\]where $\widehat f$ is the chosen normalized first arrow.
\lean{KIP126.Literature.Standard.normalized_cofiber}
% Blueprint note: Source verification (2026-10-18): Partial source correspondence. The type quantifies over every selected TriangleData and its fixed first lift. Source Notation 3.19 and Proposition 3.20 do not prove that full generality or selected-lift recognition. Primary references: [LWX126](https://arxiv.org/abs/2412.10879): Notation 3.19; Proposition 3.20; pages 19; Repository-frozen AimPaper/main.tex and AimPaper/2412.10879.pdf
% Blueprint note: Detailed source attribution: (Adapted from \cite{LWX126}, Notation 3.19 and Proposition 3.20; this general selected-map comparison remains to be verified.)
\end{definition}

\begin{definition}[Suspension and left tensoring]
\label{frontier:Literature:320f401ce5e1}
\level{fine}
(Adapted from \cite{May01}, Axiom TC1.)\par{}

For each synthetic object $X$, equip $X\otimes-$ with coherent natural isomorphisms commuting with every integral suspension.
\lean{KIP126.Literature.Standard.may_leftShift}
% Blueprint note: Source verification (2026-10-18): Source adaptation checked. These choices realize integral suspension comparisons for the existing fixed synthetic tensor; TC1–TC2, not Lemma 4.6, are the relevant framework. Compatibility of this concrete choice remains to be constructed. Primary references: [May01](https://www.math.uchicago.edu/~may/PAPERS/AddJan01.pdf): Axioms TC1–TC2; Remark 4.2; pages 11; Author-hosted January 2001 preprint, 31 pages; numbered statements retain published numbering [Pst](https://arxiv.org/pdf/1803.01804v3): Proposition 4.2; Remark 4.5; Definition 4.6; pages 38–39; arXiv:1803.01804v3, 10 November 2022
% Blueprint note: Detailed source attribution: (Adapted from \cite{May01}, Axiom TC1 and Remark 4.2, in the synthetic category of \cite{Pst}, Proposition 4.2.)
\end{definition}

\begin{definition}[Suspension and right tensoring]
\label{frontier:Literature:4983add560f1}
\level{fine}
(Adapted from \cite{May01}, Axiom TC1.)\par{}

For each synthetic object $X$, equip $-\otimes X$ with coherent natural isomorphisms commuting with every integral suspension.
\lean{KIP126.Literature.Standard.may_rightShift}
% Blueprint note: Source verification (2026-10-18): Source adaptation checked. These choices realize integral suspension comparisons for the existing fixed synthetic tensor; TC1–TC2, not Lemma 4.6, are the relevant framework. Compatibility of this concrete choice remains to be constructed. Primary references: [May01](https://www.math.uchicago.edu/~may/PAPERS/AddJan01.pdf): Axioms TC1–TC2; Remark 4.2; pages 11; Author-hosted January 2001 preprint, 31 pages; numbered statements retain published numbering [Pst](https://arxiv.org/pdf/1803.01804v3): Proposition 4.2; Remark 4.5; Definition 4.6; pages 38–39; arXiv:1803.01804v3, 10 November 2022
% Blueprint note: Detailed source attribution: (Adapted from \cite{May01}, Axiom TC1 and Remark 4.2, in the synthetic category of \cite{Pst}, Proposition 4.2.)
\end{definition}

\begin{definition}[Realization of synthetic spheres]
\label{frontier:Literature:572b7f00f725}
\level{fine}
(Adapted from \cite{Pst}, Prop. 4.40.)\par{}

For all integers $m,w$, choose sphere isomorphisms\[S^m\cong\operatorname{Re}(S^{m,w}).\]
\lean{KIP126.Literature.Standard.realization_sphere}
% Blueprint note: Source verification (2026-10-18): Source adaptation checked. The sphere coordinates, unit equality, and realization of every lambda power are proved in Lean from the fixed adjunction, shifts, lambda-localization, and bi-shift coherence. Their underlying project model constructions retain their existing proof obligations. Primary references: [Pst](https://arxiv.org/pdf/1803.01804v3): Definition 4.6; Proposition 4.40; Remark 4.41; pages 39, 49; arXiv:1803.01804v3, 10 November 2022
% Blueprint note: Detailed source attribution: (Adapted from \cite{Pst}, Definition 4.6, Proposition 4.40, and Remark 4.41.)
\end{definition}

\begin{definition}[The tmf algebra]
\label{frontier:Literature:5894b2586cbc}
\level{fine}
(Adapted from \cite{BR21}, Sec. 0.1.)\par{}

Equip $\mathit{tmf}$ with its algebra structure in the classical monoidal category.
\lean{KIP126.Literature.Standard.tmf_algebra}
% Blueprint note: Source verification (2026-10-18): Source adaptation checked. The cited papers concern the usual classical spectrum or its 2-primary/2-completed version. Recognition of the fixed project category, Adams tower, coefficient object, and selected Lean constants is a separate, currently pending formal comparison; this audit does not prove it. Primary references: [BR21](https://www.sas.rochester.edu/mth/sites/doug-ravenel/otherpapers/rognes-bruner-tmf-book.pdf): Section0.1,p1;Section1.2,p64; Author preliminary version of AMS Mathematical Surveys and Monographs 253 (2021), downloaded 2026-10-18; printed page numbers cited, PDF index = printed page + 17 in the body.
% Blueprint note: Detailed source attribution: (Adapted from the $E_\infty$ ring structure in \cite{BR21}, Section 0.1, p. 1, and Section 1.2, p. 64.)
\end{definition}

\begin{definition}[The classical w class]
\label{frontier:Literature:75fd0d147c9c}
\level{fine}
(Adapted from \cite{BR21}, Thm. 11.59.)\par{}

Fix $w\in\pi_{45}S^0$.
\lean{KIP126.Literature.Standard.wClass}
% Blueprint note: Source verification (2026-10-18): Source adaptation checked. The cited papers concern the usual classical spectrum or its 2-primary/2-completed version. Recognition of the fixed project category, Adams tower, coefficient object, and selected Lean constants is a separate, currently pending formal comparison; this audit does not prove it. Primary references: [BR21](https://www.sas.rochester.edu/mth/sites/doug-ravenel/otherpapers/rognes-bruner-tmf-book.pdf): BR21 Theorem11.59,pp457–458;Table11.3,p458,roww in(stem45,filtration9); IWX Table4,p215,row(45,9,24); BMQ Section1,pp3–4,names following Figure1.1; Author preliminary version of AMS Mathematical Surveys and Monographs 253 (2021), downloaded 2026-10-18; printed page numbers cited, PDF index = printed page + 17 in the body. [IWX](https://www.numdam.org/item/10.1007/s10240-023-00139-1.pdf): BR21 Theorem11.59,pp457–458;Table11.3,p458,roww in(stem45,filtration9); IWX Table4,p215,row(45,9,24); BMQ Section1,pp3–4,names following Figure1.1; Published IHES137 (2023),107–243; PDF index=printed page−106. [tmf](https://arxiv.org/pdf/2011.08956v4): BR21 Theorem11.59,pp457–458;Table11.3,p458,roww in(stem45,filtration9); IWX Table4,p215,row(45,9,24); BMQ Section1,pp3–4,names following Figure1.1; arXiv:2011.08956v4, 20 November 2021; published Geom. Topol.27 (2023)2763–2831. Locators below use v4 page numbering.
% Blueprint note: Detailed source attribution: (Adapted from \cite{BR21}, Theorem 11.59 and Table 11.3, pp. 457--458; notation as in \cite{IWX}, Table 4.)
\end{definition}

\begin{definition}[The spectrum of topological modular forms]
\label{frontier:Literature:79d2c0d8349e}
\level{fine}
(Project choice; cf.\ \cite{BR21}, Sec. 0.1.)\par{}

Let $\mathit{tmf}$ denote the classical spectrum of topological modular forms.
\lean{KIP126.Literature.Standard.tmf}
% Blueprint note: Source verification (2026-10-18): Project choice or construction. The cited papers concern the usual classical spectrum or its 2-primary/2-completed version. Recognition of the fixed project category, Adams tower, coefficient object, and selected Lean constants is a separate, currently pending formal comparison; this audit does not prove it. Primary references: [BR21](https://www.sas.rochester.edu/mth/sites/doug-ravenel/otherpapers/rognes-bruner-tmf-book.pdf): Section0.1,pp1–2; Author preliminary version of AMS Mathematical Surveys and Monographs 253 (2021), downloaded 2026-10-18; printed page numbers cited, PDF index = printed page + 17 in the body.
% Blueprint note: Detailed source attribution: (Project choice of the connective spectrum described in \cite{BR21}, Section 0.1, pp. 1--2.)
\end{definition}

\begin{definition}[The classical theta-five class]
\label{frontier:Literature:7bbc7d9fc81d}
\level{fine}
(Project choice; cf.\ \cite{Xu}, Cor. 1.3.)\par{}

Fix $\theta_5\in\pi_{62}S^0$.
\lean{KIP126.Literature.Standard.theta5}
% Blueprint note: Source verification (2026-10-18): Project choice or construction. This is a project choice of a 2-primary Kervaire representative. Its detection and order-two properties are separate correlated inputs, not part of this bare constant type. This verifies the mathematical source and the stated adaptation, not the pending comparison with the fixed project model. No Lean proof obligation is discharged by this audit. Primary references: [Xu](https://msp.org/gt/2016/20-3/gt-v20-n3-p08-p.pdf): Corollary 1.3; pages 1612 [BJMtheta5](https://www.sas.rochester.edu/mth/sites/doug-ravenel/otherpapers/barjoma.pdf): Theorem 2.1 and its immediate conclusion; pages [536, 537]
% Blueprint note: Detailed source attribution: (Adapted from \cite{Xu}, Corollary 1.3, p.~1612: a project choice of its order-two representative; see also \cite{BJMtheta5}, Theorem 2.1.)
\end{definition}

\begin{definition}[A synthetic h-zero class]
\label{frontier:Literature:85883e4bcc3f}
\level{fine}
(Project choice; cf.\ \cite{LWX126}, Example 5.8(4).)\par{}

Fix $[h_0]\in\pi_{0,1}S^{0,0}$.
\lean{KIP126.Literature.Standard.syntheticH0}
% Blueprint note: Source verification (2026-10-18): Project choice or construction. This entry chooses the h-zero class only. Its normalization, products and comparison claims are separate bound statements; choosing it is not itself a quoted theorem. Primary references: [LWX126](https://arxiv.org/abs/2412.10879): Theorem 7.3, Remarks 7.4–7.5; Example 5.8(4); proof of Lemma 7.16; pages 42, 28, 49–50; Repository AimPaper/main.tex and AimPaper/2412.10879.pdf snapshot; no unstated version substitution [BHS](https://arxiv.org/pdf/1910.14116v3): Remark 9.5; Theorem A.8; Corollary A.9(1); Notation A.10, Corollary A.11; Proposition A.20, relations (0), (9), and Remark A.22; pages 32, 72–73, 78–79; published Acta page mapping supplied in bhs_pages/proposals.json; arXiv:1910.14116v3, as independently retrieved by the BHS audit worker
% Blueprint note: Detailed source attribution: (Project choice following \cite{LWX126}, Example 5.8(4), and \cite{BHS}, Proposition A.20.)
\end{definition}

\begin{definition}[Universal coefficients for mod-two cohomology]
\label{frontier:Literature:87014c07c847}
\level{fine}
(Adapted from \cite{EKMM}, Thm. IV.4.5.)\par{}

For every classical spectrum $X$ and every integer $n$, there is an additive equivalence\[H^n(X;\mathbb F_2)\cong\operatorname{Hom}_{\mathrm{Ab}}(H_{-n}(X;\mathbb F_2),\mathbb F_2).\]
\lean{KIP126.Literature.Standard.universal_coefficient}
% Blueprint note: Source verification (2026-10-18): Source adaptation checked. The displayed H^n is legacy [Sigma^n X,HF2], equivalent to ordinary H^{-n}. There is no finite-dimensional or finite-type restriction. Hom_Ab coincides with F2-linear Hom because the only scalars are zero and one. This verifies the mathematical source and the stated adaptation, not the pending comparison with the fixed project model. No Lean proof obligation is discharged by this audit. Primary references: [SchwedeSymmetricSpectra](https://www.math.uni-bonn.de/~schwede/SymSpec-v3.pdf): Proposition II.6.28 and Equation (6.27); pages [281] [EKMM](https://www.math.uchicago.edu/~may/BOOKS/EKMM.pdf): Theorem IV.4.5, specialized R=N=HF2; pages [82]
% Blueprint note: Detailed source attribution: (Adapted from \cite{EKMM}, Theorem IV.4.5, specialized to $H\mathbb F_2$; equivalently \cite{SchwedeSymmetricSpectra}, Proposition II.6.28 and Equation (6.27), with the project indexing $q=-n$.)
\end{definition}

\begin{definition}[The distinguished BX theta-five]
\label{frontier:Literature:9e21a71352a4}
\level{fine}
(Project choice; cf.\ \cite{BurklundXu}, Prop. 7.19.)\par{}

Fix the distinguished synthetic class $\theta_5^{\mathrm{BX}}\in\pi_{62,64}S^{0,0}$.
\lean{KIP126.Literature.Standard.bxTheta5}
% Blueprint note: Source verification (2026-10-18): Project choice or construction. This is the project-selected BX source choice, not a claim for every class with the same label. BX uses bidegree (stem, Adams filtration); the project uses (stem, weight), with weight = stem + filtration. Thus the source theta in degree (62,2) is displayed here in degree (62,64), and the source total boundary target is regraded to the project shift (1,-1). The finite statement retains r >= 1. Source-to-model identification and the registered Lean proof obligations remain separate from this citation check. Primary references: [BurklundXu](https://arxiv.org/pdf/2302.11869v3): Proposition 7.19 and its proof; Theorem 7.20 and its proof; pages 54; arXiv:2302.11869v3, 9 November 2024; accepted version [LWX126](https://arxiv.org/abs/2412.10879): Theorem 7.3, Remarks 7.4–7.5; Example 5.8(4); proof of Lemma 7.16; pages 42, 28, 49–50; Repository AimPaper/main.tex and AimPaper/2412.10879.pdf snapshot; no unstated version substitution
% Blueprint note: Detailed source attribution: (Project choice following \cite{BurklundXu}, Proposition 7.19 and its proof.)
\end{definition}

\begin{definition}[The lambda-inverting comparison]
\label{frontier:Literature:a1cdf0a80ba8}
\level{fine}
(Adapted from \cite{BHS}, Thm. A.8.)\par{}

For every $X$ in the chosen classical family, choose a morphism of cycle-and-boundary data from the synthetic Adams system of $\nu X$ to the classical Adams system of $X$, pulled back along $(s,t,w)\mapsto(s,t)$. In particular it induces maps on all pages and on the infinite page.
\lean{KIP126.Literature.Standard.inversion_data}
% Blueprint note: Source verification (2026-10-18): Source adaptation checked. The chosen SSDataMorphism includes all cycle and boundary levels, including the project infinity construction; this is packaging beyond the literal printed spectral-sequence morphism. Write the BHS coordinates as (s,k,omega) and the project coordinates as (s,t,w). The translation is k=t-s and omega=w-(t-s); tau is renamed lambda. Hence s-omega=t-w, and the BHS differential degree (r,-1,1) becomes (r,r-1,0). This verifies the mathematical source and its stated regrading, not the outstanding Lean comparison with the fixed project model. Existing Lean proof and transitive-dependency status must remain unchanged. Spectral-sequence page and E-infinity subquotient statements here do not assert convergence to untruncated homotopy groups; BHS Convention A.5 allows an arbitrary spectrum for Theorem A.8 and Corollaries A.9/A.11. Primary references: [Pst](https://arxiv.org/pdf/1803.01804v3): Proposition 4.36; Theorem 4.37; Proposition 4.40; pages arXiv PDF pp. 47-49; published pp. 621-624; arXiv:1803.01804v3, 10 November 2022; numbering also checked in Invent. Math. 232 (2023) [BHS](https://arxiv.org/pdf/1910.14116v3): paragraph immediately before Notation A.7; proof of Theorem A.8; pages arXiv PDF p. 72; published p. 307; arXiv:1910.14116v3, 10 May 2022; checked against Acta Math. 231 (2023), 205-344 [BHS](https://arxiv.org/pdf/1910.14116v3): Theorem A.8; Notation A.6; pages arXiv PDF pp. 71-72; published pp. 306-307; arXiv:1910.14116v3, 10 May 2022; checked against Acta Math. 231 (2023), 205-344
% Blueprint note: Detailed source attribution: (Adapted from \cite{Pst}, Proposition 4.40, and \cite{BHS}, Notation A.6, the comparison before Notation A.7, and Theorem A.8.)
\end{definition}

\begin{definition}[The Ext class Delta h-one g]
\label{frontier:Literature:b367d2906c16}
\level{fine}
(Adapted from \cite{BR21}, Prop. 1.4.)\par{}

Fix $\Delta h_1g\in E_2^{9,54}(S^0)$.
\lean{KIP126.Literature.Standard.deltaH1g}
% Blueprint note: Source verification (2026-10-18): Source adaptation checked. The cited papers concern the usual classical spectrum or its 2-primary/2-completed version. Recognition of the fixed project category, Adams tower, coefficient object, and selected Lean constants is a separate, currently pending formal comparison; this audit does not prove it. Primary references: [BR21](https://www.sas.rochester.edu/mth/sites/doug-ravenel/otherpapers/rognes-bruner-tmf-book.pdf): BR21 Proposition1.4;Table1.1,p56,row(stem45,filtration9), classical Ext namew,cocycle9_20;Remark1.5,pp56–57;Figure1.10,p57; IWX Section1.5,pp117–118;Table4,p215,row(45,9,24),τΔh1g,alternate namew; Author preliminary version of AMS Mathematical Surveys and Monographs 253 (2021), downloaded 2026-10-18; printed page numbers cited, PDF index = printed page + 17 in the body. [IWX](https://www.numdam.org/item/10.1007/s10240-023-00139-1.pdf): BR21 Proposition1.4;Table1.1,p56,row(stem45,filtration9), classical Ext namew,cocycle9_20;Remark1.5,pp56–57;Figure1.10,p57; IWX Section1.5,pp117–118;Table4,p215,row(45,9,24),τΔh1g,alternate namew; Published IHES137 (2023),107–243; PDF index=printed page−106.
% Blueprint note: Detailed source attribution: (Adapted from \cite{BR21}, Proposition 1.4 and Table 1.1, pp. 56--57; notation as in \cite{IWX}, Table 4, p. 215.)
\end{definition}

\begin{definition}[The tmf target class]
\label{frontier:Literature:b9dce73bafd4}
\level{fine}
(Adapted from \cite{BR21}, Thm. 5.18.)\par{}

Fix the tmf target class $\beta^5g\in E_2^{19,114}(\mathit{tmf})$.
\lean{KIP126.Literature.Standard.tmfBetaGFour}
% Blueprint note: Source verification (2026-10-18): Source adaptation checked. The cited papers concern the usual classical spectrum or its 2-primary/2-completed version. Recognition of the fixed project category, Adams tower, coefficient object, and selected Lean constants is a separate, currently pending formal comparison; this audit does not prove it. Primary references: [BR21](https://www.sas.rochester.edu/mth/sites/doug-ravenel/otherpapers/rognes-bruner-tmf-book.pdf): BR21 Theorem5.18;Table5.4,p196;Remark5.9,p190;Table3.5,p151;Theorem1.20 proof,p77; AimPaper Appendix,paragraph immediately beforeTable5; Author preliminary version of AMS Mathematical Surveys and Monographs 253 (2021), downloaded 2026-10-18; printed page numbers cited, PDF index = printed page + 17 in the body. [LWX126](https://arxiv.org/abs/2412.10879): BR21 Theorem5.18;Table5.4,p196;Remark5.9,p190;Table3.5,p151;Theorem1.20 proof,p77; AimPaper Appendix,paragraph immediately beforeTable5; Repository AimPaper/main.tex and2412.10879.pdf, frozen baseline. Used only as secondary notation/recall; original BR21 and BMQ checked independently.
% Blueprint note: Detailed source attribution: (Adapted from \cite{BR21}, Theorem 5.18 and Table 5.4: $\beta^5g=\beta g^4$ by Table 3.5.)
\end{definition}

\begin{definition}[The free-lambda page equivalence]
\label{frontier:Literature:c8c44dc9bb6c}
\level{fine}
(Adapted from \cite{BHS}, Thm. A.8.)\par{}

For $X$ in the chosen classical family, $r\ge2$, and $r-2\le t-w$, choose an isomorphism\[E_r^{s,t,w}(\nu X)\cong E_r^{s,t}(X).\]
\lean{KIP126.Literature.Standard.free_component}
% Blueprint note: Source verification (2026-10-18): Source adaptation checked. Put a=t-w. In the polynomial-lambda spectral sequence of BHS Theorem A.8, an incoming classical d_i has a source label with lambda exponent a-i+1. All incoming differentials contributing before page r have 2<=i<=r-1, so they all occur once a>=r-2. Outgoing cycles agree under localization. Thus the localization map is an isomorphism for a>=r-2, including the boundary case; this is a derived range, not the literal range printed in Corollary A.9(4). The inverse, E2 normalization, lambda square and differential square are obtained from this same localization map, rather than from unrelated choices of isomorphism. Write the BHS coordinates as (s,k,omega) and the project coordinates as (s,t,w). The translation is k=t-s and omega=w-(t-s); tau is renamed lambda. Hence s-omega=t-w, and the BHS differential degree (r,-1,1) becomes (r,r-1,0). This verifies the mathematical source and its stated regrading, not the outstanding Lean comparison with the fixed project model. Existing Lean proof and transitive-dependency status must remain unchanged. Spectral-sequence page and E-infinity subquotient statements here do not assert convergence to untruncated homotopy groups; BHS Convention A.5 allows an arbitrary spectrum for Theorem A.8 and Corollaries A.9/A.11. Primary references: [BHS](https://arxiv.org/pdf/1910.14116v3): Theorem A.8; Notation A.6; pages arXiv PDF pp. 71-72; published pp. 306-307; arXiv:1910.14116v3, 10 May 2022; checked against Acta Math. 231 (2023), 205-344 [BHS](https://arxiv.org/pdf/1910.14116v3): Corollary A.9(1)-(4), including its final surjectivity sentence; pages arXiv PDF pp. 72-73; published pp. 307-308; arXiv:1910.14116v3, 10 May 2022; checked against Acta Math. 231 (2023), 205-344 [BHS](https://arxiv.org/pdf/1910.14116v3): paragraph immediately before Notation A.7; proof of Theorem A.8; pages arXiv PDF p. 72; published p. 307; arXiv:1910.14116v3, 10 May 2022; checked against Acta Math. 231 (2023), 205-344 [LWX126](https://arxiv.org/html/2412.10879v2#S3): Theorem 3.6 and Remark 3.7; pages local PDF p. 16; arXiv:2412.10879v2, 22 February 2025; local AimPaper/main.tex and PDF also checked
% Blueprint note: Detailed source attribution: (Adapted from \cite{BHS}, Theorem A.8 and Corollary A.9; grading as in \cite{LWX126}, Theorem 3.6 and Remark 3.7.)
\end{definition}



\begin{definition}[The finite-quotient E-two page]
\label{frontier:Literature:e30343a1c6cc}
\level{fine}
(Adapted from \cite{BHS}, Cor. A.11.)\par{}

For $X$ in the chosen classical family and $q>0$, choose isomorphisms\[E_2^{s,t,w}(\nu X/\lambda^q)\cong\begin{cases}E_2^{s,t}(X),&0\le t-w<q,\\0,&\text{otherwise}.\end{cases}\]
\lean{KIP126.Literature.Standard.finite_quotient_component}
% Blueprint note: Source verification (2026-10-18): Source adaptation checked. The canonical nu-E Adams resolution can be tensored with C-tau^q. Its polynomial E1 terms truncate to exponents 0,...,q-1; the d1 calculation then gives the stated E2 strip. At weight exponent a=t-w, passing to page r inspects outgoing d_i for i<=r-1 whose target exponent is a+i-1. The condition a+r-2<q retains all these targets and all incoming sources, so the actual quotient projection is an isomorphism on that page. This is a deduction from the tower used in BHS Theorem A.8 and Corollary A.11, not a separately numbered safe-range theorem in that paper. Write the BHS coordinates as (s,k,omega) and the project coordinates as (s,t,w). The translation is k=t-s and omega=w-(t-s); tau is renamed lambda. Hence s-omega=t-w, and the BHS differential degree (r,-1,1) becomes (r,r-1,0). This verifies the mathematical source and its stated regrading, not the outstanding Lean comparison with the fixed project model. Existing Lean proof and transitive-dependency status must remain unchanged. Spectral-sequence page and E-infinity subquotient statements here do not assert convergence to untruncated homotopy groups; BHS Convention A.5 allows an arbitrary spectrum for Theorem A.8 and Corollaries A.9/A.11. Primary references: [BHS](https://arxiv.org/pdf/1910.14116v3): Theorem A.8; Notation A.6; pages arXiv PDF pp. 71-72; published pp. 306-307; arXiv:1910.14116v3, 10 May 2022; checked against Acta Math. 231 (2023), 205-344 [BHS](https://arxiv.org/pdf/1910.14116v3): Corollary A.11(1)-(3) and its proof; Notation A.10; pages arXiv PDF p. 73; published p. 308; arXiv:1910.14116v3, 10 May 2022; checked against Acta Math. 231 (2023), 205-344 [LWX126](https://arxiv.org/html/2412.10879v2#S3): Theorem 3.6 and Remark 3.7; pages local PDF p. 16; arXiv:2412.10879v2, 22 February 2025; local AimPaper/main.tex and PDF also checked
% Blueprint note: Detailed source attribution: (Adapted from \cite{BHS}, Theorem A.8, Notation A.10 and Corollary A.11; grading as in \cite{LWX126}, Remark 3.7.)
\end{definition}

\begin{definition}[The tmf detector comparison]
\label{frontier:Literature:e51ecaf2ffbc}
\level{fine}
(Project choice; cf.\ \cite{BR21}, Sec. 0.1.)\par{}

Fix an isomorphism $\mathit{tmf}\cong D$ to the classical detector.
\lean{KIP126.Literature.Standard.tmfIso}
% Blueprint note: Source verification (2026-10-18): Project choice or construction. The cited papers concern the usual classical spectrum or its 2-primary/2-completed version. Recognition of the fixed project category, Adams tower, coefficient object, and selected Lean constants is a separate, currently pending formal comparison; this audit does not prove it. Primary references: [BR21](https://www.sas.rochester.edu/mth/sites/doug-ravenel/otherpapers/rognes-bruner-tmf-book.pdf): Section0.1,pp1–2 (source object only); Author preliminary version of AMS Mathematical Surveys and Monographs 253 (2021), downloaded 2026-10-18; printed page numbers cited, PDF index = printed page + 17 in the body.
% Blueprint note: Detailed source attribution: (Project comparison with the spectrum described in \cite{BR21}, Section 0.1; this is not a theorem of the cited source.)
\end{definition}

\begin{theorem}[Filtration of the sphere action]
\label{frontier:Literature:3dfd6c694616}
\level{fine}
(Adapted from \cite{Moss}, Sec. 1; \cite{BHS}, Remark 9.5.)\par{}

Let $X$ lie in the chosen synthetic object family. For integers $m,w,m^{\prime},w^{\prime},s,s^{\prime}$, take $a\in\pi_{m,w}S^{0,0}$ and $b\in\pi_{m^{\prime},w^{\prime}}X$. If $a$ has Adams filtration at least $s$ and $b$ has filtration at least $s^{\prime}$, then $a\cdot b$ has filtration at least $s+s^{\prime}$.
\lean{KIP126.Literature.Standard.action_filtration}
% Blueprint note: Source verification (2026-10-18): Partial source correspondence. The three bindings are general comparison statements for the project-selected Adams towers, not a single numbered theorem in AimPaper. Moss Section 1 supplies composition pairings in classical Adams systems, while BHS Remark 9.5 and the canonical towers in Appendix A explain the synthetic monoidal setting. The exact compatibility of the selected finite-quotient E2 coordinates, action and abutment for every positive quotient length, and the action-filtration statement on all shifts and iterated quotients, must still be supplied. No upper bound is inserted into the frozen formal range; this citation check does not silently narrow it to the paper examples. Primary references: [Moss](https://gdz.sub.uni-goettingen.de/download/pdf/PPN266833020_0115/LOG_0047.pdf): Section 1; Theorem (1.2) and conditions (1.3); pages 283–285; Original Math. Z. 115 (1970), 283–310; GDZ digitized journal scan LOG_0047 [BHS](https://arxiv.org/pdf/1910.14116v3): Remark 9.5; Theorem A.8; Corollary A.9(1); Notation A.10, Corollary A.11; Proposition A.20, relations (0), (9), and Remark A.22; pages 32, 72–73, 78–79; published Acta page mapping supplied in bhs_pages/proposals.json; arXiv:1910.14116v3, as independently retrieved by the BHS audit worker [LWX126](https://arxiv.org/abs/2412.10879): Theorem 7.3, Remarks 7.4–7.5; Example 5.8(4); proof of Lemma 7.16; pages 42, 28, 49–50; Repository AimPaper/main.tex and AimPaper/2412.10879.pdf snapshot; no unstated version substitution
% Blueprint note: Detailed source attribution: (Adapted from the multiplicative framework of \cite{Moss}, Section 1, and \cite{BHS}, Remark 9.5 and Appendix A; the selected filtration and finite-quotient comparisons remain explicit obligations.)
\end{theorem}

\begin{theorem}[Detection of the finite-quotient action]
\label{frontier:Literature:5ea1d02075a4}
\level{fine}
(Adapted from \cite{Moss}, Sec. 1; \cite{BHS}, Remark 9.5.)\par{}

Let $q>0$ and $s,t,s^{\prime},t^{\prime},k,l\in\mathbb N$. Suppose $x\in E_2^{s,t}(S^0)$ detects $a\in\pi_{t-s,t-k}S^{0,0}$ under the synthetic comparison, and the projected label of $y\in E_2^{s^{\prime},t^{\prime}}(S^0)$ detects $b\in\pi_{t^{\prime}-s^{\prime},t^{\prime}-l}(S^{0,0}/\lambda^q)$. Then the projected label of $xy$ detects $a\cdot b$ in tridegree $(s+s^{\prime},t+t^{\prime},t+t^{\prime}-k-l)$.
\lean{KIP126.Literature.Standard.finite_action}
% Blueprint note: Source verification (2026-10-18): Partial source correspondence. The three bindings are general comparison statements for the project-selected Adams towers, not a single numbered theorem in AimPaper. Moss Section 1 supplies composition pairings in classical Adams systems, while BHS Remark 9.5 and the canonical towers in Appendix A explain the synthetic monoidal setting. The exact compatibility of the selected finite-quotient E2 coordinates, action and abutment for every positive quotient length, and the action-filtration statement on all shifts and iterated quotients, must still be supplied. No upper bound is inserted into the frozen formal range; this citation check does not silently narrow it to the paper examples. Primary references: [Moss](https://gdz.sub.uni-goettingen.de/download/pdf/PPN266833020_0115/LOG_0047.pdf): Section 1; Theorem (1.2) and conditions (1.3); pages 283–285; Original Math. Z. 115 (1970), 283–310; GDZ digitized journal scan LOG_0047 [BHS](https://arxiv.org/pdf/1910.14116v3): Remark 9.5; Theorem A.8; Corollary A.9(1); Notation A.10, Corollary A.11; Proposition A.20, relations (0), (9), and Remark A.22; pages 32, 72–73, 78–79; published Acta page mapping supplied in bhs_pages/proposals.json; arXiv:1910.14116v3, as independently retrieved by the BHS audit worker [LWX126](https://arxiv.org/abs/2412.10879): Theorem 7.3, Remarks 7.4–7.5; Example 5.8(4); proof of Lemma 7.16; pages 42, 28, 49–50; Repository AimPaper/main.tex and AimPaper/2412.10879.pdf snapshot; no unstated version substitution
% Blueprint note: Detailed source attribution: (Adapted from the multiplicative framework of \cite{Moss}, Section 1, and \cite{BHS}, Remark 9.5 and Appendix A; the selected filtration and finite-quotient comparisons remain explicit obligations.)
\end{theorem}

\begin{theorem}[Classical products and the tmf Hurewicz map]
\label{frontier:Literature:860e3947f28a}
\level{fine}
(Adapted from \cite{BR21}, Thm. 11.61(20); \cite{tmf}, Thm. 1.2.)\par{}

\begin{enumerate}
\item Let $s,t,s^{\prime},t^{\prime}\in\mathbb N$. If $x\in E_2^{s,t}(S^0)$ detects $a\in\pi_{t-s}S^0$ and $y\in E_2^{s^{\prime},t^{\prime}}(S^0)$ detects $b\in\pi_{t^{\prime}-s^{\prime}}S^0$, then $xy$ detects $ab$ in bidegree $(s+s^{\prime},t+t^{\prime})$.
\item The isomorphism $\mathit{tmf}\cong D$ carries the tmf unit to the detector unit.
\item The class $g$ detects $\bar\kappa$.
\item For $s\le0$, $E_2^{s,63+s}(\mathit{tmf})=0$.
\item No positive odd integer annihilates the tmf Hurewicz image of $\bar\kappa^4w$.
\item $\pi_{62}\mathit{tmf}=0$.
\item The class $\Delta h_1g$ detects $w$.
\end{enumerate}
\lean{KIP126.Literature.Standard.classical_product_detection,KIP126.Literature.Standard.tmfUnit_agrees,KIP126.Literature.Standard.kappaBar_detection,KIP126.Literature.Standard.tmf_low_filtration63,KIP126.Literature.Standard.tmf_high125_nonzero,KIP126.Literature.Standard.tmf_vanishing62,KIP126.Literature.Standard.w_detection}
% Blueprint note: Source verification (2026-10-18): Partial source correspondence. The cited papers concern the usual classical spectrum or its 2-primary/2-completed version. Recognition of the fixed project category, Adams tower, coefficient object, and selected Lean constants is a separate, currently pending formal comparison; this audit does not prove it. Primary references: [BR21](https://www.sas.rochester.edu/mth/sites/doug-ravenel/otherpapers/rognes-bruner-tmf-book.pdf): Section0.3,pp5–6 (multiplicative stronglyconvergentAdamssequence);Section1.2,p64 (algebra spectral sequence); Section0.1,p1 (source ring unit only); Theorem11.61,case(20),p467;Table11.3,p458; Section1.2,p64,E2=Ext_A(2)(F2,F2);Section0.3,p5;Section0.5,pp8–9; BMQ Figure1.1,p3,nonzeroκbar^4w atstem125;Section1,pp3–4;Theorem1.2,p4;discussionpp5–6; BR21 Theorem9.27,pp338–340;Table9.4,stem125row; BR21 Theorem11.59,pp457–458;Table11.3,p458,roww; IWX Table4,p215,row(45,9,24); Author preliminary version of AMS Mathematical Surveys and Monographs 253 (2021), downloaded 2026-10-18; printed page numbers cited, PDF index = printed page + 17 in the body. [tmf](https://arxiv.org/pdf/2011.08956v4): BMQ Figure1.1,p3,nonzeroκbar^4w atstem125;Section1,pp3–4;Theorem1.2,p4;discussionpp5–6; BR21 Theorem9.27,pp338–340;Table9.4,stem125row; Figure1.1,p3;chart conventions in Section1,pp2–4; arXiv:2011.08956v4, 20 November 2021; published Geom. Topol.27 (2023)2763–2831. Locators below use v4 page numbering. [IWX](https://www.numdam.org/item/10.1007/s10240-023-00139-1.pdf): BR21 Theorem11.59,pp457–458;Table11.3,p458,roww; IWX Table4,p215,row(45,9,24); Published IHES137 (2023),107–243; PDF index=printed page−106.
% Blueprint note: Detailed source attribution: (Adapted from \cite{BR21}, Section 0.3, Theorem 11.61(20), and Table 11.3, and \cite{tmf}, Figure 1.1 and Theorem 1.2; the detector-unit comparison is project data.)
\end{theorem}





\begin{theorem}[Quotient multiplication and the sphere action]
\label{frontier:Literature:b4eb6770b658}
\level{fine}
(Adapted from \cite{LWX126}, Remark 3.4.)\par{}

Let $q>0$, $x\in\pi_{m,n}S^{0,0}$ and $y\in\pi_{k,l}(S^{0,0}/\lambda^q)$, with all four indices integral. Multiplication of the quotient image of $x$ by $y$ equals the sphere action $x\cdot y$.
\lean{KIP126.Literature.Standard.quotientAlgebra_sphere_action}
% Blueprint note: Source verification (2026-10-18): Source adaptation checked. The equality is the unital algebra/module compatibility of sphere action with quotient multiplication. Its exact formula depends on the project’s tensor-to-shift and unit identifications and remains a separate comparison obligation. Primary references: [LWX126](https://arxiv.org/abs/2412.10879): Remark 3.4; pages 15; Repository-frozen AimPaper/main.tex and AimPaper/2412.10879.pdf [Pst](https://arxiv.org/pdf/1803.01804v3): Definition 4.6; pages 39; arXiv:1803.01804v3, 10 November 2022
% Blueprint note: Detailed source attribution: (Adapted from the quotient algebra in \cite{LWX126}, Remark 3.4, and the sphere tensor grading in \cite{Pst}, Definition 4.6; the selected-model action comparison remains to be verified.)
\end{theorem}





\begin{theorem}[Realization of lambda and the unit]
\label{frontier:Literature:00344e96581e}
\level{fine}
(Adapted from \cite{Pst}, Prop. 4.40.)\par{}

For all $m,w\in\mathbb Z$ and $k\in\mathbb N$, realization carries $\lambda^k:S^{m,w-k}\to S^{m,w}$ to the identity in sphere coordinates. At $(m,w)=(0,0)$, the sphere isomorphism agrees with the composite of the zero-shift identification, the inverse $\nu$--realization isomorphism, and the realization of the synthetic unit isomorphism.
\lean{KIP126.Literature.Standard.realization_lambda,KIP126.Literature.Standard.realization_unit}
% Blueprint note: Source verification (2026-10-18): Source adaptation checked. The sphere coordinates, unit equality, and realization of every lambda power are proved in Lean from the fixed adjunction, shifts, lambda-localization, and bi-shift coherence. Their underlying project model constructions retain their existing proof obligations. Primary references: [Pst](https://arxiv.org/pdf/1803.01804v3): Definition 4.6; Proposition 4.40; Remark 4.41; pages 39, 49; arXiv:1803.01804v3, 10 November 2022
% Blueprint note: Detailed source attribution: (Adapted from \cite{Pst}, Proposition 4.40 and Remark 4.41; using the specified unit and lambda comparisons.)
\end{theorem}

\begin{theorem}[Nonzero permanence of the theta-five label]
\label{frontier:Literature:bc100fe38326}
\level{fine}
(Adapted from \cite{BJMtheta5}, Thm. 2.1.)\par{}

The class $h_5^2\in E_2^{2,64}(S^0)$ is a nonzero permanent cycle.
\lean{KIP126.Literature.Standard.h5Square_permanent}
% Blueprint note: Source verification (2026-10-18): Source adaptation checked. The selected theta5 must be the source detection witness. First-quotient or fixed-model comparison is not silently supplied by the citation. This verifies the mathematical source and the stated adaptation, not the pending comparison with the fixed project model. No Lean proof obligation is discharged by this audit. Primary references: [BJMtheta5](https://www.sas.rochester.edu/mth/sites/doug-ravenel/otherpapers/barjoma.pdf): Theorem 2.1 and its immediate conclusion; pages [536, 537] [Xu](https://msp.org/gt/2016/20-3/gt-v20-n3-p08-p.pdf): Corollary 1.3; pages 1612 [IWXCharts](https://zenodo.org/records/6987157/files/Adams-classical-Einfty.csv?download=1): Adams-classical-Einfty.csv, physical rows 204--207 (including header) [IWX](https://arxiv.org/pdf/2001.04511v3): Corollary 1.2, with the cited classical Adams charts; pages 4
% Blueprint note: Detailed source attribution: (Adapted from \cite{BJMtheta5}, Theorem 2.1, pp.~536--537; \cite{Xu}, Corollary 1.3, p.~1612; see also \cite{IWX}, Corollary 1.2.)
\end{theorem}

\begin{theorem}[Detection of the chosen theta-five class]
\label{frontier:Literature:71f4461cf400}
\level{fine}
(Adapted from \cite{BJMtheta5}, Thm. 2.1.)\par{}

The standard class $h_5^2\in E_2^{2,64}(S^0)$ detects the chosen $\theta_5\in\pi_{62}S^0$.
\lean{KIP126.Literature.Standard.theta5_detection}
% Blueprint note: Source verification (2026-10-18): Source adaptation checked. The selected theta5 must be the source detection witness. First-quotient or fixed-model comparison is not silently supplied by the citation. This verifies the mathematical source and the stated adaptation, not the pending comparison with the fixed project model. No Lean proof obligation is discharged by this audit. Primary references: [BJMtheta5](https://www.sas.rochester.edu/mth/sites/doug-ravenel/otherpapers/barjoma.pdf): Theorem 2.1 and its immediate conclusion; pages [536, 537] [Xu](https://msp.org/gt/2016/20-3/gt-v20-n3-p08-p.pdf): Corollary 1.3; pages 1612 [IWXCharts](https://zenodo.org/records/6987157/files/Adams-classical-Einfty.csv?download=1): Adams-classical-Einfty.csv, physical rows 204--207 (including header) [IWX](https://arxiv.org/pdf/2001.04511v3): Corollary 1.2, with the cited classical Adams charts; pages 4
% Blueprint note: Detailed source attribution: (Adapted from \cite{BJMtheta5}, Theorem 2.1, pp.~536--537; \cite{Xu}, Corollary 1.3, p.~1612; see also \cite{IWX}, Corollary 1.2.)
\end{theorem}

\begin{theorem}[Normalized nu and retained synthetic labels]
\label{frontier:Literature:0385573fba72}
\level{fine}
(Adapted from \cite{LWX126}, Example 6.19; \cite{BHS}, Thm. A.8.)\par{}

\begin{enumerate}
\item The normalized Hopf map $\nu\in\pi_{3,4}S^{0,0}$ has first-quotient label $h_2$.
\item Let $q>0$ and $0\le t-w<q$. Projection to $\nu X/\lambda^q$, followed by the finite-quotient $E_2$ coordinates, sends the weight-$w$ lift of each $x\in E_2^{s,t}(X)$ to $x$.
\item For the Hopf-$\nu$ cofiber triangle, whenever its three normalized exponents sum to one, its three normalized arrows form a distinguished triangle.
\item For every $X$ in the chosen classical family and $r\ge2$, $E_r^{s,t,w}(\nu X)=0$ when $t<w$.
\end{enumerate}
\lean{KIP126.Literature.Standard.nuCofiber_normalized_label,KIP126.Literature.Standard.finite_quotient_component_retained,KIP126.Literature.Standard.nuCofiber_triangle,KIP126.Literature.Standard.above_diagonal}
% Blueprint note: Source verification (2026-10-18): Partial source correspondence. This mixed entry contains four independent claims. The page statements are verified after regrading BHS (s,k,omega) to project (s,t,w), where k=t-s and omega=w-k. The Hopf-nu first label and distinguishedness of the three already selected normalized arrows require explicit object/arrow comparisons; existence of some normalized triangle does not identify arbitrary chosen lifts. Primary references: [LWX126](https://arxiv.org/abs/2412.10879): Theorem 7.3, Remarks 7.4–7.5; Example 5.8(4); proof of Lemma 7.16; pages 42, 28, 49–50; Repository AimPaper/main.tex and AimPaper/2412.10879.pdf snapshot; no unstated version substitution [BHS](https://arxiv.org/pdf/1910.14116v3): Remark 9.5; Theorem A.8; Corollary A.9(1); Notation A.10, Corollary A.11; Proposition A.20, relations (0), (9), and Remark A.22; pages 32, 72–73, 78–79; published Acta page mapping supplied in bhs_pages/proposals.json; arXiv:1910.14116v3, as independently retrieved by the BHS audit worker
% Blueprint note: Detailed source attribution: (Adapted from \cite{LWX126}, Example 6.19 and Proposition 3.20, and \cite{BHS}, Theorem A.8, Corollary A.9(1), Notation A.10 and Corollary A.11. The selected-map comparisons remain explicit obligations.)
\end{theorem}

\begin{theorem}[Exactness of tensoring]
\label{frontier:Literature:1a607d05c555}
\level{fine}
(Adapted from \cite{May01}, Axiom TC2.)\par{}

For every synthetic object $X$, both $X\otimes-$ and $-\otimes X$, equipped with their chosen integral shift isomorphisms, preserve distinguished triangles.
\lean{KIP126.Literature.Standard.may_leftExact,KIP126.Literature.Standard.may_rightExact}
% Blueprint note: Source verification (2026-10-18): Source adaptation checked. Colimit preservation implies exactness for the synthetic stable category; the Lean declarations require this for the exact suspension structures selected above. Primary references: [May01](https://www.math.uchicago.edu/~may/PAPERS/AddJan01.pdf): Axioms TC1–TC2; Remark 4.2; pages 11; Author-hosted January 2001 preprint, 31 pages; numbered statements retain published numbering [Pst](https://arxiv.org/pdf/1803.01804v3): Proposition 4.2; Remark 4.5; Definition 4.6; pages 38–39; arXiv:1803.01804v3, 10 November 2022
% Blueprint note: Detailed source attribution: (Adapted from \cite{May01}, Axiom TC2 and Remark 4.2, in the synthetic category of \cite{Pst}, Proposition 4.2.)
\end{theorem}

\begin{theorem}[Infinite representatives in the untruncated window]
\label{frontier:Literature:322d0d20f79e}
\level{fine}
(Adapted from \cite{BHS}, Cor. A.9.)\par{}

Let $X$ belong to the chosen classical family, $w\le t$, $x\in E_2^{s,t}(X)$ and $z\in E_\infty^{s,t,w}(\nu X)$. The weight-$w$ label of $x$ represents $z$ if and only if $x\in Z_\infty^{s,t}(X)$ and the window image of $z$ is the class of $x$ modulo $B_{1+t-w}^{s,t}(X)$.
\lean{KIP126.Literature.Standard.eInfinity_nu_labels}
% Blueprint note: Source verification (2026-10-18): Source adaptation checked. BHS Corollary A.9 gives the untruncated E-infinity subquotient, and Corollary A.11 gives the finite quotient. All formula indices transform by s-omega=t-w. The finite formula holds for q>=1 in BHS, while AimPaper Proposition 3.12 explicitly states q>=2. Additive equivalences are viewed as Z-linear without a change of scalar structure. Matching a specific E2 representative to its quotient class requires the canonical label-compatible presentation, not an arbitrary abstract isomorphism. Zero labels are permitted in subgroup and quotient statements. Write the BHS coordinates as (s,k,omega) and the project coordinates as (s,t,w). The translation is k=t-s and omega=w-(t-s); tau is renamed lambda. Hence s-omega=t-w, and the BHS differential degree (r,-1,1) becomes (r,r-1,0). This verifies the mathematical source and its stated regrading, not the outstanding Lean comparison with the fixed project model. Existing Lean proof and transitive-dependency status must remain unchanged. Spectral-sequence page and E-infinity subquotient statements here do not assert convergence to untruncated homotopy groups; BHS Convention A.5 allows an arbitrary spectrum for Theorem A.8 and Corollaries A.9/A.11. Primary references: [BHS](https://arxiv.org/pdf/1910.14116v3): Corollary A.9(1)-(4), including its final surjectivity sentence; pages arXiv PDF pp. 72-73; published pp. 307-308; arXiv:1910.14116v3, 10 May 2022; checked against Acta Math. 231 (2023), 205-344 [LWX126](https://arxiv.org/html/2412.10879v2#S3): Proposition 3.11; pages local PDF p. 17; arXiv:2412.10879v2, 22 February 2025; local AimPaper/main.tex and PDF also checked
% Blueprint note: Detailed source attribution: (Adapted from \cite{BHS}, Corollary A.9, and \cite{LWX126}, Proposition 3.11.)
\end{theorem}

\begin{theorem}[Finite lifts and reaching a page]
\label{frontier:Literature:3aad75009ef0}
\level{fine}
(Adapted from \cite{BHS}, Thm. A.1(1a)--(1b).)\par{}

Let $X$ belong to the chosen classical family, $q>0$, and $x\in E_2^{s,t}(X)$. The label $x$ reaches $E_{q+1}$ if and only if there is $a\in\pi_{t-s,t}(\nu X/\lambda^q)$ whose restriction to $\nu X/\lambda$ is the first-quotient class of $x$. Reaching a page permits zero and boundaries.
\lean{KIP126.Literature.Standard.bhs_finite_lift}
% Blueprint note: Source verification (2026-10-18): Source adaptation checked. Finite-quotient result: write q for BHS r, identify tau with lambda, and use stem k=t-s. Reaching E_(q+1) means d_2 through d_q vanish; zero and earlier boundaries are allowed. Unlike the untruncated conclusion, the finite case extends to arbitrary spectra via Corollary A.11 and Lemma A.14, as in the proof of Theorem A.1(1). Over F_2 the sign in A.1(1c) disappears. The mathematical source has been checked; identification with the project's specified coefficient object, Adams towers, quotient maps and realization remains part of the existing formal comparison obligations. No Lean proof obligation is discharged by this citation. Primary references: [BHS](https://arxiv.org/pdf/1910.14116v3): Theorem A.1(1a)–(1c); Corollary A.11; Lemma A.14; proof of Theorem 9.19(1); pages PDF pp. 70, 73–75; arXiv:1910.14116v3, 10 May 2022; final prepublication version, not published pagination [LWX126](https://arxiv.org/abs/2412.10879): Theorem 3.8(1)–(2); pages PDF p. 16; Repository AimPaper/main.tex and AimPaper/2412.10879.pdf; exact local source snapshot
% Blueprint note: Detailed source attribution: (Adapted from \cite{BHS}, Theorem A.1(1a)--(1b), Corollary A.11 and Lemma A.14; \cite{LWX126}, Theorem 3.8.)
\end{theorem}

\begin{theorem}[Adams filtration and lambda divisibility]
\label{frontier:Literature:436802be2502}
\level{fine}
(Adapted from \cite{BHS}, Cor. A.19.)\par{}

For $X$ in the chosen classical family, integers $m,w,s$ with $w-m\le s$, and $a\in\pi_{m,w}\nu X$, the class $a$ has Adams filtration at least $s$ if and only if there is $b\in\pi_{m,m+s}\nu X$ with $\lambda^{m+s-w}b=a$.
\lean{KIP126.Literature.Standard.selected_filtration_lambda}
% Blueprint note: Source verification (2026-10-18): Partial source correspondence. BHS proves the source result for E-nilpotent complete spectra with a strongly convergent Adams spectral sequence. The selected family includes the ordinary sphere, its nu-cofiber, the detector, and shifts; membership alone does not establish BHS completeness. This entry is a pending source-to-selected-model comparison, not a direct application of BHS. The mathematical source has been checked; identification with the project's specified coefficient object, Adams towers, quotient maps and realization remains part of the existing formal comparison obligations. No Lean proof obligation is discharged by this citation. Primary references: [BHS](https://arxiv.org/pdf/1910.14116v3): Corollary A.19; pages PDF p. 76; arXiv:1910.14116v3, 10 May 2022; final prepublication version, not published pagination
% Blueprint note: Detailed source attribution: (Adapted from \cite{BHS}, Corollary A.19. Source-to-selected-model comparison pending.)
\end{theorem}

\begin{theorem}[The first-quotient label of eta]
\label{frontier:Literature:48d4b249a919}
\level{fine}
(Adapted from \cite{LWX126}, Notation 7.2.)\par{}

The first-quotient image of the normalized synthetic class $\eta\in\pi_{1,2}S^{0,0}$ is $h_1\in E_2^{1,2}(S^0)$.
\lean{KIP126.Literature.Standard.eta_label}
% Blueprint note: Source verification (2026-10-18): Source adaptation checked. Notation 7.2 explicitly uses eta=[h1] in pi_(1,2); Theorem 3.8(1) identifies the first quotient with E2. The chosen project eta must satisfy this comparison, which remains a registered proof obligation. Primary references: [LWX126](https://arxiv.org/abs/2412.10879): Theorem 7.3, Remarks 7.4–7.5; Example 5.8(4); proof of Lemma 7.16; pages 42, 28, 49–50; Repository AimPaper/main.tex and AimPaper/2412.10879.pdf snapshot; no unstated version substitution
% Blueprint note: Detailed source attribution: (Adapted from \cite{LWX126}, Notation 7.2, using the first-quotient convention of Theorem 3.8(1).)
\end{theorem}

\begin{theorem}[Restriction between finite infinite-page windows]
\label{frontier:Literature:636e4c89b3e4}
\level{fine}
(Adapted from \cite{LWX126}, Remark 3.14.)\par{}

Let $X$ belong to the chosen classical family, $0<i\le j$, and $0\le t-w<i$. Under the infinite-page window equivalences, restriction $\nu X/\lambda^j\to\nu X/\lambda^i$ induces the natural map\[Z_{j-t+w}^{s,t}/B_{1+t-w}^{s,t}\longrightarrow Z_{i-t+w}^{s,t}/B_{1+t-w}^{s,t}.\]
\lean{KIP126.Literature.Standard.eInfinity_rho_finite}
% Blueprint note: Source verification (2026-10-18): Source adaptation checked. BHS Corollary A.9 gives the untruncated E-infinity subquotient, and Corollary A.11 gives the finite quotient. All formula indices transform by s-omega=t-w. The finite formula holds for q>=1 in BHS, while AimPaper Proposition 3.12 explicitly states q>=2. Additive equivalences are viewed as Z-linear without a change of scalar structure. Matching a specific E2 representative to its quotient class requires the canonical label-compatible presentation, not an arbitrary abstract isomorphism. Zero labels are permitted in subgroup and quotient statements. Write the BHS coordinates as (s,k,omega) and the project coordinates as (s,t,w). The translation is k=t-s and omega=w-(t-s); tau is renamed lambda. Hence s-omega=t-w, and the BHS differential degree (r,-1,1) becomes (r,r-1,0). This verifies the mathematical source and its stated regrading, not the outstanding Lean comparison with the fixed project model. Existing Lean proof and transitive-dependency status must remain unchanged. Spectral-sequence page and E-infinity subquotient statements here do not assert convergence to untruncated homotopy groups; BHS Convention A.5 allows an arbitrary spectrum for Theorem A.8 and Corollaries A.9/A.11. Primary references: [BHS](https://arxiv.org/pdf/1910.14116v3): Corollary A.11(1)-(3) and its proof; Notation A.10; pages arXiv PDF p. 73; published p. 308; arXiv:1910.14116v3, 10 May 2022; checked against Acta Math. 231 (2023), 205-344 [BHS](https://arxiv.org/pdf/1910.14116v3): Corollary A.9(1)-(4), including its final surjectivity sentence; pages arXiv PDF pp. 72-73; published pp. 307-308; arXiv:1910.14116v3, 10 May 2022; checked against Acta Math. 231 (2023), 205-344 [LWX126](https://arxiv.org/html/2412.10879v2#S3): Remark 3.14; pages local PDF p. 17; arXiv:2412.10879v2, 22 February 2025; local AimPaper/main.tex and PDF also checked
% Blueprint note: Detailed source attribution: (Adapted from \cite{BHS}, Corollary A.11 and its proof, and \cite{LWX126}, Remark 3.14.)
\end{theorem}

\begin{theorem}[Detection of the classical Hopf nu]
\label{frontier:Literature:64496b8f7295}
\level{fine}
(Adapted from \cite{Adams58}, Lem. 2.3.)\par{}

The label $h_2\in E_2^{1,4}(S^0)$ detects the classical Hopf map $\nu\in\pi_3S^0$.
\lean{KIP126.Literature.Standard.nu_detection}
% Blueprint note: Source verification (2026-10-18): Source adaptation checked. This fixes the actual classical Hopf nu and the standard h2 label on the project Adams sequence. This verifies the mathematical source and the stated adaptation, not the pending comparison with the fixed project model. No Lean proof obligation is discharged by this audit. Primary references: [Adams58](https://www.sas.rochester.edu/mth/sites/doug-ravenel/otherpapers/Adams-SS.pdf): Lemma 2.3 and the following m=0 paragraph; pages 184 [Adams60](https://www.maths.ed.ac.uk/~v1ranick/papers/adams1.pdf): Section 1.1, after Theorem 1.1.1; pages 20
% Blueprint note: Detailed source attribution: (Adapted from \cite{Adams58}, Lemma 2.3, p.~184; and \cite{Adams60}, Section 1.1, p.~20.)
\end{theorem}

\begin{theorem}[The tmf power-operation differential]
\label{frontier:Literature:6c63a4b40657}
\level{fine}
(Adapted from \cite{BR21}, Thm. 5.18.)\par{}

The tmf labels in bidegrees $(16,112)$ and $(19,114)$ satisfy\[d_3(v_2^{16})=\beta^5g.\]
\lean{KIP126.Literature.Standard.tmf_br21}
% Blueprint note: Source verification (2026-10-18): Source adaptation checked. The cited papers concern the usual classical spectrum or its 2-primary/2-completed version. Recognition of the fixed project category, Adams tower, coefficient object, and selected Lean constants is a separate, currently pending formal comparison; this audit does not prove it. Primary references: [BR21](https://www.sas.rochester.edu/mth/sites/doug-ravenel/otherpapers/rognes-bruner-tmf-book.pdf): Theorem5.18;Table5.4,p196;Proposition5.8(4)andRemark5.9,p190;Table3.5,p151;Table1.3,p65; Author preliminary version of AMS Mathematical Surveys and Monographs 253 (2021), downloaded 2026-10-18; printed page numbers cited, PDF index = printed page + 17 in the body.
% Blueprint note: Detailed source attribution: (Adapted from \cite{BR21}, Theorem 5.18 and Table 5.4, using $\beta^4=g^3$ from Table 3.5.)
\end{theorem}

\begin{theorem}[Weight shift preserves the E-two label]
\label{frontier:Literature:6d7a8d6a798c}
\level{fine}
(Adapted from \cite{BHS}, Thm. A.8(2).)\par{}

For $X$ in the chosen classical family, $s,t\in\mathbb Z$, $k\in\mathbb N$ and $x\in E_2^{s,t}(X)$, weight shift\[E_2^{s,t,t-k}(\nu X)\cong E_2^{s,t,t-k-1}(\Sigma^{0,-1}\nu X)\]carries the unshifted synthetic label of $x$ to its label on $\Sigma^{0,-1}\nu X$.
\lean{KIP126.Literature.Standard.nuPage_e2}
% Blueprint note: Source verification (2026-10-18): Source adaptation checked. The public shift isomorphism includes integer pages below 2; those pages are frozen to the first page by a project convention, not covered as extra pages in the cited paper. Write the BHS coordinates as (s,k,omega) and the project coordinates as (s,t,w). The translation is k=t-s and omega=w-(t-s); tau is renamed lambda. Hence s-omega=t-w, and the BHS differential degree (r,-1,1) becomes (r,r-1,0). This verifies the mathematical source and its stated regrading, not the outstanding Lean comparison with the fixed project model. Existing Lean proof and transitive-dependency status must remain unchanged. Spectral-sequence page and E-infinity subquotient statements here do not assert convergence to untruncated homotopy groups; BHS Convention A.5 allows an arbitrary spectrum for Theorem A.8 and Corollaries A.9/A.11. Primary references: [Pst](https://arxiv.org/pdf/1803.01804v3): Definition 4.6; Definition 4.27; pages arXiv PDF pp. 39, 45; published pp. 609, 618; arXiv:1803.01804v3, 10 November 2022; numbering also checked in Invent. Math. 232 (2023) [BHS](https://arxiv.org/pdf/1910.14116v3): Theorem A.8; Notation A.6; pages arXiv PDF pp. 71-72; published pp. 306-307; arXiv:1910.14116v3, 10 May 2022; checked against Acta Math. 231 (2023), 205-344 [LWX126](https://arxiv.org/html/2412.10879v2#S3): Theorem 3.6 and Remark 3.7; pages local PDF p. 16; arXiv:2412.10879v2, 22 February 2025; local AimPaper/main.tex and PDF also checked
% Blueprint note: Detailed source attribution: (Adapted from \cite{Pst}, Definition 4.6, and \cite{BHS}, Theorem A.8(2); grading as in \cite{LWX126}, Theorem 3.6 and Remark 3.7.)
\end{theorem}

\begin{theorem}[The Hopf-invariant-one differentials]
\label{frontier:Literature:75582be1124c}
\level{fine}
(From \cite{LWX126}, Introduction.)\par{}

For every $j\in\mathbb N$ with $j\ge4$, there is a nonzero differential\[d_2(h_j)=h_0h_{j-1}^2\in E_2^{3,1+2^j}(S^0).\]
\lean{KIP126.Literature.Standard.adamsOneLine_d2}
% Blueprint note: Source verification (2026-10-18): Source adaptation checked. The differential convention is d_r:(s,t)->(s+r,t+r-1). For s=1,r=2,t=2^j the target is (3,2^j+1). This verifies the mathematical source and the stated adaptation, not the pending comparison with the fixed project model. No Lean proof obligation is discharged by this audit. Primary references: [LWX126](https://arxiv.org/abs/2412.10879): Introduction, unnumbered displayed Hopf-invariant-one differential; pages 2
% Blueprint note: Detailed source attribution: (From \cite{LWX126}, Introduction, p.~2, the displayed Hopf-invariant-one differential, citing \cite{Adams60}.)
\end{theorem}

\begin{theorem}[The total boundary of h-six square]
\label{frontier:Literature:86d633a93ab1}
\level{fine}
(Adapted from \cite{BurklundXu}, proof of Prop. 7.19.)\par{}

The total boundary of the first-quotient class labeled $h_6^2$ is $\lambda\eta(\theta_5^{\mathrm{BX}})^2$.
\lean{KIP126.Literature.Standard.bx_total_boundary}
% Blueprint note: Source verification (2026-10-18): Source adaptation checked. This is the project-selected BX source choice, not a claim for every class with the same label. BX uses bidegree (stem, Adams filtration); the project uses (stem, weight), with weight = stem + filtration. Thus the source theta in degree (62,2) is displayed here in degree (62,64), and the source total boundary target is regraded to the project shift (1,-1). The finite statement retains r >= 1. Source-to-model identification and the registered Lean proof obligations remain separate from this citation check. Primary references: [BurklundXu](https://arxiv.org/pdf/2302.11869v3): Proposition 7.19 and its proof; Theorem 7.20 and its proof; pages 54; arXiv:2302.11869v3, 9 November 2024; accepted version [LWX126](https://arxiv.org/abs/2412.10879): Theorem 7.3, Remarks 7.4–7.5; Example 5.8(4); proof of Lemma 7.16; pages 42, 28, 49–50; Repository AimPaper/main.tex and AimPaper/2412.10879.pdf snapshot; no unstated version substitution
% Blueprint note: Detailed source attribution: (Adapted from \cite{BurklundXu}, proof of Proposition 7.19; regraded as in \cite{LWX126}, Remark 7.5.)
\end{theorem}

\begin{theorem}[A torsion lift of a boundary]
\label{frontier:Literature:90a18b9fe1e0}
\level{fine}
(Adapted from \cite{BHS}, Thm. A.1(3a).)\par{}

Let $X$ belong to the chosen classical family, $r\ge2$, and $y\in E_2^{s,t}(X)$. If $y$ has a common infinite-cycle representative and is hit on page $r$, there is $a\in\pi_{t-s,t}\nu X$ with first-quotient label $y$ and $\lambda^{r-1}a=0$.
\lean{KIP126.Literature.Standard.selected_realization_boundary_lift}
% Blueprint note: Source verification (2026-10-18): Partial source correspondence. BHS proves the source result for E-nilpotent complete spectra with a strongly convergent Adams spectral sequence. The selected family includes the ordinary sphere, its nu-cofiber, the detector, and shifts; membership alone does not establish BHS completeness. This entry is a pending source-to-selected-model comparison, not a direct application of BHS. The mathematical source has been checked; identification with the project's specified coefficient object, Adams towers, quotient maps and realization remains part of the existing formal comparison obligations. No Lean proof obligation is discharged by this citation. Primary references: [BHS](https://arxiv.org/pdf/1910.14116v3): Theorem A.1(3a), equivalently Theorem 9.19(3a); pages PDF pp. 70, 77 (also p. 36); arXiv:1910.14116v3, 10 May 2022; final prepublication version, not published pagination
% Blueprint note: Detailed source attribution: (Adapted from \cite{BHS}, Theorem A.1(3a). Source-to-selected-model comparison pending.)
\end{theorem}

\begin{theorem}[Classical inputs and synthetic comparison identities]
\label{frontier:Literature:d19d6b7e62ef}
\level{fine}
(Sources are cited individually below.)\par{}

Use the complete-source hypotheses of Definition~\ref{sem:Literature:bhs-applicability}. All unqualified synthetic comparisons below concern the chosen classical object family.
\begin{enumerate}
\item (Adapted from \cite{Xu}, Cor. 1.3.)\par{}

$2\theta_5=0$.
\item (Adapted from \cite{BHS}, Thm. A.8(2).)\par{}

For $X$ in the chosen classical family and $w\le t$, the lambda-inverting comparison sends the weight-$w$ label of each $x\in E_2^{s,t}(X)$ to $x$.
\item (Adapted from \cite{SchwedeSymmetricSpectra}, Prop. II.5.22.)\par{}

If a classical spectrum $X$ is $n$-connected, then $H_i(X;\mathbb F_2)=0$ for every $i\le n$.
\item (Adapted from \cite{Adams58}, Lem. 2.3.)\par{}

The label $h_2\in E_2^{1,4}(S^0)$ is a nonzero permanent cycle.
\item (Adapted from \cite{BHS}, Thm. A.8(2).)\par{}

For $X$ in the chosen classical family and $w\le t$, the free-range isomorphism on $E_2$ sends the weight-$w$ label of $x\in E_2^{s,t}(X)$ to $x$.
\item (Adapted from \cite{BHS}, Thm. A.1(2a).)\par{}

Assume the complete-source hypotheses on $X$. If $r\ge1$, $x\in E_2^{s,t}(X)$ has a common permanent representative and survives nontrivially to $E_{r+1}$, and $a\in\pi_{t-s,t}\nu X$ has first-quotient label $x$, then $\lambda^{r-1}a\ne0$.
\item (Adapted from \cite{BHS}, Thm. A.1(2)--(3).)\par{}

The preceding lifetime conclusion holds for every $X$ in the chosen classical family, with the same hypotheses on $r,x,a$.
\item (Adapted from \cite{BHS}, Lem. 9.15.)\par{}

For any classical map $f:X\to Y$ of Adams tower filtration at least $k\in\mathbb N$, there is a map $\widetilde f:\nu X\to\Sigma^{0,-k}\nu Y$ with $\lambda^k\widetilde f=\nu f$.
\item (Adapted from \cite{BHS}, Thm. A.1(2b), (3b).)\par{}

Assume the complete-source hypotheses on $X$. If a nonzero permanent label $x\in E_2^{s,t}(X)$ detects a prescribed $a\in\pi_{t-s}X$, then $a$ is the realization of a class in $\pi_{t-s,t}\nu X$ with first-quotient label $x$.
\item (Adapted from \cite{BHS}, Thm. A.1(3a).)\par{}

Assume the complete-source hypotheses on $X$. If $r\ge2$ and $y\in E_2^{s,t}(X)$ has a common infinite-cycle representative and is hit on page $r$, it has a lift in $\pi_{t-s,t}\nu X$ annihilated by $\lambda^{r-1}$.
\item (Adapted from \cite{BHS}, Thm. A.8.)\par{}

For $r\ge2$ and $r-2\le t-w$, the free-range isomorphism commutes with $d_r$, with target bidegree $(s+r,t+r-1)$.
\item (Adapted from \cite{BHS}, Thm. A.8.)\par{}

For $r\ge2$ and $r-2\le t-w$, the free-range isomorphism equals the lambda-inverting comparison.
\item (Adapted from \cite{BR21}, Prop. 1.4.)\par{}

Every nonzero class in $E_2^{4,24}(S^0)$ equals $g$.
\item (Adapted from \cite{BHS}, Thm. A.1(2b), (3b).)\par{}

Assume the complete-source hypotheses on $X$. If $x\in E_2^{s,t}(X)$ is a nonzero permanent label and $a\in\pi_{t-s,t}\nu X$ has first-quotient label $x$, then $x$ detects the realization of $a$.
\item (Adapted from \cite{Moss}, Sec. 1; \cite{BelmontKong2024}, Def. 2.3.)\par{}

For objects $X,Y$ in the chosen mapping family, let $T$ be their mapping Adams tower. If $f\in\pi_n T_{s+1}$ lifts to every later stage $T_t$ with $t\ge s+1$ and maps to zero in $\pi_n T_s$, then $f=0$, for every $s\in\mathbb N$ and $n\in\mathbb Z$.
\item (Adapted from \cite{Adams58}, Lem. 2.3.)\par{}

The label $h_1\in E_2^{1,2}(S^0)$ is a nonzero permanent cycle.
\item (Adapted from \cite{BR21}, Sec. 0.3.)\par{}

For every integer $n$, $H_n(\mathit{tmf};\mathbb F_2)$ is finite.
\item (Adapted from \cite{SchwedeSymmetricSpectra}, Example I.1.8.)\par{}

The classical sphere is $(-1)$-connected.
\item (Adapted from \cite{BHS}, Thm. A.8.)\par{}

For every $X$ in the chosen classical family and $r\ge2$, the lambda-inverting comparison commutes with $d_r$ in all weights.
\item (Adapted from \cite{BHS}, Thm. A.8.)\par{}

For $r\ge2$ and $r-2\le t-w$, multiplication by $\lambda$ followed by the free-range comparison at weight $w-1$ equals the free-range comparison at weight $w$.
\item (Adapted from \cite{BHS}, Cor. A.19.)\par{}

Assume the complete-source hypotheses on $X$. A label $x\in E_2^{s,t}(X)$ has a common infinite-cycle representative if and only if it is the first-quotient image of an untruncated class in $\pi_{t-s,t}\nu X$.
\item (Adapted from \cite{Pst}, Lem. 4.23.)\par{}

For every distinguished classical triangle $T$, its image under $\nu$ with the same first two arrows admits a distinguished synthetic triangle if and only if the mod-two homology sequence of $T$ is short exact.
\item (Adapted from \cite{BHS}, Thm. A.1(2)--(3).)\par{}

For every $X$ in the chosen classical family, a nonzero permanent label $x\in E_2^{s,t}(X)$ detects the realization of every $a\in\pi_{t-s,t}\nu X$ with first-quotient label $x$.
\item (Adapted from \cite{BR21}, Prop. 1.4.)\par{}

Every nonzero class in $E_2^{9,54}(S^0)$ equals $\Delta h_1g$.
\item (Adapted from \cite{LWX126}, Prop. 3.16.)\par{}

The realization of the normalized synthetic $\eta\in\pi_{1,2}S^{0,0}$ is the classical Hopf map $\eta$.
\item (Adapted from \cite{BHS}, Thm. A.8(3).)\par{}

For every $X$ in the chosen classical family and $N\in\mathbb Z$, degeneration of its classical Adams spectral sequence at $N$ implies degeneration of that of $\nu X$ at $N$.
\item (Adapted from \cite{LWX126}, Prop. 3.17.)\par{}

Let $T$ be a distinguished classical triangle with mod-two homology short exact and connecting map of Adams filtration at least one. There is a distinguished synthetic lift of $T$ whose connecting arrow $\delta$ satisfies $\lambda\delta=\nu(T.h)$. For each full lift of $T.h$ through $\lambda^k$ with $k\ge1$, its reduction to one lambda division differs from $\delta$ by lambda torsion.
\item (Adapted from \cite{Adams58}, Lem. 2.3.)\par{}

The label $h_0\in E_2^{1,1}(S^0)$ detects the classical degree-two map.
\item (Adapted from \cite{BR21}, Prop. 1.4.)\par{}

$g\ne0$.
\item (From \cite{LWX126}, proof of Lem. 7.10.)\par{}

If $a,b\in\pi_{62}S^0$ are both detected by $h_5^2$, their difference has Adams filtration at least six.
\item (Adapted from \cite{BHS}, Cor. A.19.)\par{}

Assume the complete-source hypotheses on $X$. For $w-m\le s$ and $a\in\pi_{m,w}\nu X$, Adams filtration at least $s$ is equivalent to the existence of $b\in\pi_{m,m+s}\nu X$ with $\lambda^{m+s-w}b=a$.
\item (Adapted from \cite{BR21}, Sec. 0.1.)\par{}

For every $n<0$, $\pi_n\mathit{tmf}=0$.
\item (Cf.\ \cite{LWX126}, Prop. 3.20.)\par{}

If the three normalized filtration exponents of a classical triangle in the chosen family sum to one, then there exist arrows forming a distinguished triangle on $\nu X,\nu Y,\nu Z$ with connecting target $\Sigma\nu X$.
\item (Adapted from \cite{May01}, Lem. 4.6.)\par{}

For two synthetic cofiber triangles $T,U$, there is an object $V$ with maps to $T.X\otimes U.Z$, $T.Y\otimes U.Y$, and $T.Z\otimes U.X$, commuting with the two tensor-product squares. Its two boundary composites are negatives of one another. Each pair of homotopy classes in the first two targets agreeing in $T.Y\otimes U.Z$ lifts to a common class on $V$.
\item (Adapted from \cite{BHS}, Thm. A.8.)\par{}

For every $X$ in the chosen classical family and $r\ge2$, multiplication by $\lambda$ followed by the lambda-inverting comparison at weight $w-1$ equals that comparison at weight $w$.
\item (Adapted from \cite{BR21}, Prop. 1.4.)\par{}

$\Delta h_1g\ne0$.
\end{enumerate}
\lean{KIP126.Literature.Standard.theta5_order_two,KIP126.Literature.Standard.inversion_e2,KIP126.Literature.Standard.connected_homology_vanishing,KIP126.Literature.Standard.nu_permanent,KIP126.Literature.Standard.free_e2,KIP126.Literature.Standard.bhs_realization_lifetime,KIP126.Literature.Standard.selected_realization_lifetime,KIP126.Literature.Standard.lift,KIP126.Literature.Standard.bhs_realization_prescribed_lift,KIP126.Literature.Standard.bhs_realization_boundary_lift,KIP126.Literature.Standard.free_differential,KIP126.Literature.Standard.free_inversion,KIP126.Literature.Standard.g_unique,KIP126.Literature.Standard.bhs_realization_detection,KIP126.Literature.Standard.moss_weak_convergence,KIP126.Literature.Standard.eta_permanent,KIP126.Literature.Standard.tmf_finiteModTwoType,KIP126.Literature.Standard.sphere_connected,KIP126.Literature.Standard.inversion_differential,KIP126.Literature.Standard.free_lambda,KIP126.Literature.Standard.bhs_permanent_lift,KIP126.Literature.Standard.nu_cofiber,KIP126.Literature.Standard.selected_realization_detection,KIP126.Literature.Standard.deltaH1g_unique,KIP126.Literature.Standard.eta_realization,KIP126.Literature.Standard.degeneration,KIP126.Literature.Standard.triangle_lift,KIP126.Literature.Standard.two_detection,KIP126.Literature.Standard.g_nonzero,KIP126.Literature.Standard.theta5_filtration_gap,KIP126.Literature.Standard.bhs_filtration_lambda,KIP126.Literature.Standard.tmf_connective,KIP126.Literature.Standard.unshifted,KIP126.Literature.Standard.may_pushpull,KIP126.Literature.Standard.inversion_lambda,KIP126.Literature.Standard.deltaH1g_nonzero}
% Blueprint note: Source verification (2026-10-18): Partial source correspondence. Each of the 36 clauses has its own primary source and locator. Source verification does not discharge any existing Lean proof or source-to-model comparison. Clause 26 has an unverified pre-initial-page extension; clause 35 omits the weight shifts in its related source and is not verified. The other clauses retain their explicit completeness, convergence, coordinate and selected-model qualifications. Primary references: [Xu](https://msp.org/gt/2016/20-3/gt-v20-n3-p08-p.pdf): Corollary 1.3; pages 1612 [IWX](https://arxiv.org/pdf/2001.04511v3): Corollary 1.3, Table 1, row k=62; pages [4, 6] [SchwedeSymmetricSpectra](https://www.math.uni-bonn.de/~schwede/SymSpec-v3.pdf): Proposition II.5.22; Example I.1.14; Definition II.6.21; pages [16, 17, 264, 265, 278, 279] [EKMM](https://www.math.uchicago.edu/~may/BOOKS/EKMM.pdf): Proposition IV.1.2(i), after shifting X by -(n+1); pages [72] [Adams58](https://www.sas.rochester.edu/mth/sites/doug-ravenel/otherpapers/Adams-SS.pdf): Lemma 2.3 and the following m=0 paragraph; pages 184 [Adams60](https://www.maths.ed.ac.uk/~v1ranick/papers/adams1.pdf): Section 1.1, after Theorem 1.1.1; pages 20 [SchwedeSymmetricSpectra](https://www.math.uni-bonn.de/~schwede/SymSpec-v3.pdf): Example I.1.8; proof of Proposition II.5.21; pages [12, 264] [LWX126](https://arxiv.org/abs/2412.10879): Proposition 3.16; Notation 3.19; Notation 7.2; pages [18, 19, 42] [Pst](https://arxiv.org/pdf/1803.01804v3): Proposition 4.40 and Remark 4.41; pages 49 [LWX126](https://arxiv.org/abs/2412.10879): Proof of Lemma 7.10; pages 45 [IWXCharts](https://zenodo.org/records/6987157/files/Adams-classical-Einfty.csv?download=1): Adams-classical-Einfty.csv, physical rows 204--207 (including header) [BR21](https://www.sas.rochester.edu/mth/sites/doug-ravenel/otherpapers/rognes-bruner-tmf-book.pdf): Proposition1.4;Table1.1,p54;Remark1.5,pp56–57;Figure1.9,p56 [BR21](https://www.sas.rochester.edu/mth/sites/doug-ravenel/otherpapers/rognes-bruner-tmf-book.pdf): Section0.3,p5;Section0.5,pp8–9;Section1.2,p64 [BR21](https://www.sas.rochester.edu/mth/sites/doug-ravenel/otherpapers/rognes-bruner-tmf-book.pdf): BR21 Proposition1.4;Table1.1,p56;Remark1.5,pp56–57;Figure1.10,p57; IWX Table4,p215,row(45,9,24) [IWX](https://www.numdam.org/item/10.1007/s10240-023-00139-1.pdf): BR21 Proposition1.4;Table1.1,p56;Remark1.5,pp56–57;Figure1.10,p57; IWX Table4,p215,row(45,9,24) [BR21](https://www.sas.rochester.edu/mth/sites/doug-ravenel/otherpapers/rognes-bruner-tmf-book.pdf): Proposition1.4;Table1.1,p54;Figure1.9,p56 [BR21](https://www.sas.rochester.edu/mth/sites/doug-ravenel/otherpapers/rognes-bruner-tmf-book.pdf): Section0.1,p1;Section0.5,pp8–9;Section1.2,p64 [BR21](https://www.sas.rochester.edu/mth/sites/doug-ravenel/otherpapers/rognes-bruner-tmf-book.pdf): BR21 Proposition1.4;Table1.1,p56;Figure1.10,p57; IWX Table4,p215,row(45,9,24) [IWX](https://www.numdam.org/item/10.1007/s10240-023-00139-1.pdf): BR21 Proposition1.4;Table1.1,p56;Figure1.10,p57; IWX Table4,p215,row(45,9,24) [Pst](https://arxiv.org/pdf/1803.01804v3): Proposition 4.36; Theorem 4.37; Proposition 4.40; pages arXiv PDF pp. 47-49; published pp. 621-624; arXiv:1803.01804v3, 10 November 2022; numbering also checked in Invent. Math. 232 (2023) [BHS](https://arxiv.org/pdf/1910.14116v3): paragraph immediately before Notation A.7; proof of Theorem A.8; pages arXiv PDF p. 72; published p. 307; arXiv:1910.14116v3, 10 May 2022; checked against Acta Math. 231 (2023), 205-344 [BHS](https://arxiv.org/pdf/1910.14116v3): Theorem A.8; Notation A.6; pages arXiv PDF pp. 71-72; published pp. 306-307; arXiv:1910.14116v3, 10 May 2022; checked against Acta Math. 231 (2023), 205-344 [BHS](https://arxiv.org/pdf/1910.14116v3): Theorem A.1(2a), equivalently Theorem 9.19(2a); pages PDF pp. 70, 77 (also p. 36); arXiv:1910.14116v3, 10 May 2022; final prepublication version, not published pagination [BHS](https://arxiv.org/pdf/1910.14116v3): Theorem A.1(2b) and (3b), equivalently Theorem 9.19(2b),(3b); pages PDF pp. 70, 77 (also p. 36); arXiv:1910.14116v3, 10 May 2022; final prepublication version, not published pagination [BHS](https://arxiv.org/pdf/1910.14116v3): Theorem A.1(3a), equivalently Theorem 9.19(3a); pages PDF pp. 70, 77 (also p. 36); arXiv:1910.14116v3, 10 May 2022; final prepublication version, not published pagination [BHS](https://arxiv.org/pdf/1910.14116v3): Theorem A.1, permanent-cycle paragraph preceding (2a); proof of Theorem 9.19(2)–(4); pages PDF pp. 70, 76–77; arXiv:1910.14116v3, 10 May 2022; final prepublication version, not published pagination [BHS](https://arxiv.org/pdf/1910.14116v3): Corollary A.19; pages PDF p. 76; arXiv:1910.14116v3, 10 May 2022; final prepublication version, not published pagination [BHS](https://arxiv.org/pdf/1910.14116v3): Lemma 9.15 and its proof; pages 34–35; arXiv:1910.14116v3, 10 May 2022 [LWX126](https://arxiv.org/abs/2412.10879): Propositions 3.16–3.17; Remark 3.18; pages 18–19; Repository-frozen AimPaper/main.tex and AimPaper/2412.10879.pdf [Pst](https://arxiv.org/pdf/1803.01804v3): Lemma 4.23; pages 43; arXiv:1803.01804v3, 10 November 2022 [LWX126](https://arxiv.org/abs/2412.10879): Notation 3.19; Proposition 3.20; pages 19; Repository-frozen AimPaper/main.tex and AimPaper/2412.10879.pdf [May01](https://www.math.uchicago.edu/~may/PAPERS/AddJan01.pdf): Axiom TC3; Lemma 4.6; pages 12–14; Author-hosted January 2001 preprint, 31 pages; numbered statements retain published numbering [Moss](https://gdz.sub.uni-goettingen.de/download/pdf/PPN266833020_0115/LOG_0047.pdf): Section 1; Theorem (1.2) and conditions (1.3); pages 283–285; Original Math. Z. 115 (1970), 283–310; GDZ digitized journal scan LOG_0047 [BelmontKong2024](https://arxiv.org/pdf/2112.08689v2): Definitions 2.3–2.4; Theorem 4.10 and Assumptions 3.14, 3.17; pages 5, 12–13, 16–17; arXiv:2112.08689v2, 16 August 2024 (not a different published numbering) [MayMatric1969](https://www.math.uchicago.edu/~may/PAPERS/8.pdf): Definition 1.2 and Proposition 2.3 (ordinary triple-product case); pages 538, 543–544; Original J. Algebra 12 (1969), 533–568, author-hosted scan
% Blueprint note: Clause 1: Selected 2-primary representative; no claim about arbitrary integral representatives. This verifies the mathematical source and the stated adaptation, not the pending comparison with the fixed project model. No Lean proof obligation is discharged by this audit. Primary references: [Xu](https://msp.org/gt/2016/20-3/gt-v20-n3-p08-p.pdf): Corollary 1.3; pages 1612 [IWX](https://arxiv.org/pdf/2001.04511v3): Corollary 1.3, Table 1, row k=62; pages [4, 6]
% Blueprint note: Clause 2: Write the BHS coordinates as (s,k,omega) and the project coordinates as (s,t,w). The translation is k=t-s and omega=w-(t-s); tau is renamed lambda. Hence s-omega=t-w, and the BHS differential degree (r,-1,1) becomes (r,r-1,0). This verifies the mathematical source and its stated regrading, not the outstanding Lean comparison with the fixed project model. Existing Lean proof and transitive-dependency status must remain unchanged. Spectral-sequence page and E-infinity subquotient statements here do not assert convergence to untruncated homotopy groups; BHS Convention A.5 allows an arbitrary spectrum for Theorem A.8 and Corollaries A.9/A.11. Primary references: [Pst](https://arxiv.org/pdf/1803.01804v3): Proposition 4.36; Theorem 4.37; Proposition 4.40; pages arXiv PDF pp. 47-49; published pp. 621-624; arXiv:1803.01804v3, 10 November 2022; numbering also checked in Invent. Math. 232 (2023) [BHS](https://arxiv.org/pdf/1910.14116v3): paragraph immediately before Notation A.7; proof of Theorem A.8; pages arXiv PDF p. 72; published p. 307; arXiv:1910.14116v3, 10 May 2022; checked against Acta Math. 231 (2023), 205-344 [BHS](https://arxiv.org/pdf/1910.14116v3): Theorem A.8; Notation A.6; pages arXiv PDF pp. 71-72; published pp. 306-307; arXiv:1910.14116v3, 10 May 2022; checked against Acta Math. 231 (2023), 205-344
% Blueprint note: Clause 3: Smash connectivity with HF2; all integer connectivity bounds. This verifies the mathematical source and the stated adaptation, not the pending comparison with the fixed project model. No Lean proof obligation is discharged by this audit. Primary references: [SchwedeSymmetricSpectra](https://www.math.uni-bonn.de/~schwede/SymSpec-v3.pdf): Proposition II.5.22; Example I.1.14; Definition II.6.21; pages [16, 17, 264, 265, 278, 279] [EKMM](https://www.math.uchicago.edu/~may/BOOKS/EKMM.pdf): Proposition IV.1.2(i), after shifting X by -(n+1); pages [72]
% Blueprint note: Clause 4: The quaternionic Hopf map supplies a nonzero surviving h2. This verifies the mathematical source and the stated adaptation, not the pending comparison with the fixed project model. No Lean proof obligation is discharged by this audit. Primary references: [Adams58](https://www.sas.rochester.edu/mth/sites/doug-ravenel/otherpapers/Adams-SS.pdf): Lemma 2.3 and the following m=0 paragraph; pages 184 [Adams60](https://www.maths.ed.ac.uk/~v1ranick/papers/adams1.pdf): Section 1.1, after Theorem 1.1.1; pages 20
% Blueprint note: Clause 5: Write the BHS coordinates as (s,k,omega) and the project coordinates as (s,t,w). The translation is k=t-s and omega=w-(t-s); tau is renamed lambda. Hence s-omega=t-w, and the BHS differential degree (r,-1,1) becomes (r,r-1,0). This verifies the mathematical source and its stated regrading, not the outstanding Lean comparison with the fixed project model. Existing Lean proof and transitive-dependency status must remain unchanged. Spectral-sequence page and E-infinity subquotient statements here do not assert convergence to untruncated homotopy groups; BHS Convention A.5 allows an arbitrary spectrum for Theorem A.8 and Corollaries A.9/A.11. Primary references: [BHS](https://arxiv.org/pdf/1910.14116v3): Theorem A.8; Notation A.6; pages arXiv PDF pp. 71-72; published pp. 306-307; arXiv:1910.14116v3, 10 May 2022; checked against Acta Math. 231 (2023), 205-344 [BHS](https://arxiv.org/pdf/1910.14116v3): paragraph immediately before Notation A.7; proof of Theorem A.8; pages arXiv PDF p. 72; published p. 307; arXiv:1910.14116v3, 10 May 2022; checked against Acta Math. 231 (2023), 205-344
% Blueprint note: Clause 6: BHS assumes E-nilpotent completeness and a strongly convergent E-based Adams spectral sequence (Definitions 9.16–9.17). BHSObjectApplicability is the project's explicit Adams-tower encoding of these hypotheses, not a quoted definition from BHS. The mathematical source has been checked; identification with the project's specified coefficient object, Adams towers, quotient maps and realization remains part of the existing formal comparison obligations. No Lean proof obligation is discharged by this citation. Primary references: [BHS](https://arxiv.org/pdf/1910.14116v3): Theorem A.1(2a), equivalently Theorem 9.19(2a); pages PDF pp. 70, 77 (also p. 36); arXiv:1910.14116v3, 10 May 2022; final prepublication version, not published pagination
% Blueprint note: Clause 7: BHS proves the source result for E-nilpotent complete spectra with a strongly convergent Adams spectral sequence. The selected family includes the ordinary sphere, its nu-cofiber, the detector, and shifts; membership alone does not establish BHS completeness. This entry is a pending source-to-selected-model comparison, not a direct application of BHS. The mathematical source has been checked; identification with the project's specified coefficient object, Adams towers, quotient maps and realization remains part of the existing formal comparison obligations. No Lean proof obligation is discharged by this citation. Primary references: [BHS](https://arxiv.org/pdf/1910.14116v3): Theorem A.1(2a), equivalently Theorem 9.19(2a); pages PDF pp. 70, 77 (also p. 36); arXiv:1910.14116v3, 10 May 2022; final prepublication version, not published pagination [BHS](https://arxiv.org/pdf/1910.14116v3): Theorem A.1(2b) and (3b), equivalently Theorem 9.19(2b),(3b); pages PDF pp. 70, 77 (also p. 36); arXiv:1910.14116v3, 10 May 2022; final prepublication version, not published pagination [BHS](https://arxiv.org/pdf/1910.14116v3): Theorem A.1(3a), equivalently Theorem 9.19(3a); pages PDF pp. 70, 77 (also p. 36); arXiv:1910.14116v3, 10 May 2022; final prepublication version, not published pagination
% Blueprint note: Clause 8: BHS factors a map of filtration k. The interface uses the actual target Adams-tower lower-bound predicate and negative-target shifts; identifying that predicate with source Adams filtration is part of model recognition. Primary references: [BHS](https://arxiv.org/pdf/1910.14116v3): Lemma 9.15 and its proof; pages 34–35; arXiv:1910.14116v3, 10 May 2022 [LWX126](https://arxiv.org/abs/2412.10879): Propositions 3.16–3.17; Remark 3.18; pages 18–19; Repository-frozen AimPaper/main.tex and AimPaper/2412.10879.pdf
% Blueprint note: Clause 9: BHS assumes E-nilpotent completeness and a strongly convergent E-based Adams spectral sequence (Definitions 9.16–9.17). BHSObjectApplicability is the project's explicit Adams-tower encoding of these hypotheses, not a quoted definition from BHS. The mathematical source has been checked; identification with the project's specified coefficient object, Adams towers, quotient maps and realization remains part of the existing formal comparison obligations. No Lean proof obligation is discharged by this citation. Primary references: [BHS](https://arxiv.org/pdf/1910.14116v3): Theorem A.1(2b) and (3b), equivalently Theorem 9.19(2b),(3b); pages PDF pp. 70, 77 (also p. 36); arXiv:1910.14116v3, 10 May 2022; final prepublication version, not published pagination
% Blueprint note: Clause 10: BHS assumes E-nilpotent completeness and a strongly convergent E-based Adams spectral sequence (Definitions 9.16–9.17). BHSObjectApplicability is the project's explicit Adams-tower encoding of these hypotheses, not a quoted definition from BHS. The mathematical source has been checked; identification with the project's specified coefficient object, Adams towers, quotient maps and realization remains part of the existing formal comparison obligations. No Lean proof obligation is discharged by this citation. Primary references: [BHS](https://arxiv.org/pdf/1910.14116v3): Theorem A.1(3a), equivalently Theorem 9.19(3a); pages PDF pp. 70, 77 (also p. 36); arXiv:1910.14116v3, 10 May 2022; final prepublication version, not published pagination
% Blueprint note: Clause 11: Write the BHS coordinates as (s,k,omega) and the project coordinates as (s,t,w). The translation is k=t-s and omega=w-(t-s); tau is renamed lambda. Hence s-omega=t-w, and the BHS differential degree (r,-1,1) becomes (r,r-1,0). This verifies the mathematical source and its stated regrading, not the outstanding Lean comparison with the fixed project model. Existing Lean proof and transitive-dependency status must remain unchanged. Spectral-sequence page and E-infinity subquotient statements here do not assert convergence to untruncated homotopy groups; BHS Convention A.5 allows an arbitrary spectrum for Theorem A.8 and Corollaries A.9/A.11. Primary references: [BHS](https://arxiv.org/pdf/1910.14116v3): Theorem A.8; Notation A.6; pages arXiv PDF pp. 71-72; published pp. 306-307; arXiv:1910.14116v3, 10 May 2022; checked against Acta Math. 231 (2023), 205-344 [BHS](https://arxiv.org/pdf/1910.14116v3): paragraph immediately before Notation A.7; proof of Theorem A.8; pages arXiv PDF p. 72; published p. 307; arXiv:1910.14116v3, 10 May 2022; checked against Acta Math. 231 (2023), 205-344
% Blueprint note: Clause 12: Write the BHS coordinates as (s,k,omega) and the project coordinates as (s,t,w). The translation is k=t-s and omega=w-(t-s); tau is renamed lambda. Hence s-omega=t-w, and the BHS differential degree (r,-1,1) becomes (r,r-1,0). This verifies the mathematical source and its stated regrading, not the outstanding Lean comparison with the fixed project model. Existing Lean proof and transitive-dependency status must remain unchanged. Spectral-sequence page and E-infinity subquotient statements here do not assert convergence to untruncated homotopy groups; BHS Convention A.5 allows an arbitrary spectrum for Theorem A.8 and Corollaries A.9/A.11. Primary references: [BHS](https://arxiv.org/pdf/1910.14116v3): Theorem A.8; Notation A.6; pages arXiv PDF pp. 71-72; published pp. 306-307; arXiv:1910.14116v3, 10 May 2022; checked against Acta Math. 231 (2023), 205-344 [BHS](https://arxiv.org/pdf/1910.14116v3): paragraph immediately before Notation A.7; proof of Theorem A.8; pages arXiv PDF p. 72; published p. 307; arXiv:1910.14116v3, 10 May 2022; checked against Acta Math. 231 (2023), 205-344
% Blueprint note: Clause 13: The cited papers concern the usual classical spectrum or its 2-primary/2-completed version. Recognition of the fixed project category, Adams tower, coefficient object, and selected Lean constants is a separate, currently pending formal comparison; this audit does not prove it. Primary references: [BR21](https://www.sas.rochester.edu/mth/sites/doug-ravenel/otherpapers/rognes-bruner-tmf-book.pdf): Proposition1.4;Table1.1,p54;Remark1.5,pp56–57;Figure1.9,p56
% Blueprint note: Clause 14: BHS assumes E-nilpotent completeness and a strongly convergent E-based Adams spectral sequence (Definitions 9.16–9.17). BHSObjectApplicability is the project's explicit Adams-tower encoding of these hypotheses, not a quoted definition from BHS. The mathematical source has been checked; identification with the project's specified coefficient object, Adams towers, quotient maps and realization remains part of the existing formal comparison obligations. No Lean proof obligation is discharged by this citation. Primary references: [BHS](https://arxiv.org/pdf/1910.14116v3): Theorem A.1(2b) and (3b), equivalently Theorem 9.19(2b),(3b); pages PDF pp. 70, 77 (also p. 36); arXiv:1910.14116v3, 10 May 2022; final prepublication version, not published pagination
% Blueprint note: Clause 15: Moss p.285 gives weak convergence for the mod-p cohomology Adams system when the source is finite and target cohomology is finitely generated degreewise; spheres satisfy those conditions. The project expresses this as residual injectivity on its chosen tower, matching Belmont--Kong Definition 2.3. Composition detection is derived here from the fixed mapping pairing and its abutment compatibility; it is not a separate conclusion of Moss Theorem 1.2. Primary references: [Moss](https://gdz.sub.uni-goettingen.de/download/pdf/PPN266833020_0115/LOG_0047.pdf): Section 1; Theorem (1.2) and conditions (1.3); pages 283–285; Original Math. Z. 115 (1970), 283–310; GDZ digitized journal scan LOG_0047 [BelmontKong2024](https://arxiv.org/pdf/2112.08689v2): Definitions 2.3–2.4; Theorem 4.10 and Assumptions 3.14, 3.17; pages 5, 12–13, 16–17; arXiv:2112.08689v2, 16 August 2024 (not a different published numbering)
% Blueprint note: Clause 16: The complex Hopf map supplies a nonzero surviving h1. This verifies the mathematical source and the stated adaptation, not the pending comparison with the fixed project model. No Lean proof obligation is discharged by this audit. Primary references: [Adams58](https://www.sas.rochester.edu/mth/sites/doug-ravenel/otherpapers/Adams-SS.pdf): Lemma 2.3 and the following m=0 paragraph; pages 184 [Adams60](https://www.maths.ed.ac.uk/~v1ranick/papers/adams1.pdf): Section 1.1, after Theorem 1.1.1; pages 20
% Blueprint note: Clause 17: The cited papers concern the usual classical spectrum or its 2-primary/2-completed version. Recognition of the fixed project category, Adams tower, coefficient object, and selected Lean constants is a separate, currently pending formal comparison; this audit does not prove it. Primary references: [BR21](https://www.sas.rochester.edu/mth/sites/doug-ravenel/otherpapers/rognes-bruner-tmf-book.pdf): Section0.3,p5;Section0.5,pp8–9;Section1.2,p64
% Blueprint note: Clause 18: Negative stable stems vanish; convert i<0 into i<=-1. This verifies the mathematical source and the stated adaptation, not the pending comparison with the fixed project model. No Lean proof obligation is discharged by this audit. Primary references: [SchwedeSymmetricSpectra](https://www.math.uni-bonn.de/~schwede/SymSpec-v3.pdf): Example I.1.8; proof of Proposition II.5.21; pages [12, 264]
% Blueprint note: Clause 19: Write the BHS coordinates as (s,k,omega) and the project coordinates as (s,t,w). The translation is k=t-s and omega=w-(t-s); tau is renamed lambda. Hence s-omega=t-w, and the BHS differential degree (r,-1,1) becomes (r,r-1,0). This verifies the mathematical source and its stated regrading, not the outstanding Lean comparison with the fixed project model. Existing Lean proof and transitive-dependency status must remain unchanged. Spectral-sequence page and E-infinity subquotient statements here do not assert convergence to untruncated homotopy groups; BHS Convention A.5 allows an arbitrary spectrum for Theorem A.8 and Corollaries A.9/A.11. Primary references: [BHS](https://arxiv.org/pdf/1910.14116v3): paragraph immediately before Notation A.7; proof of Theorem A.8; pages arXiv PDF p. 72; published p. 307; arXiv:1910.14116v3, 10 May 2022; checked against Acta Math. 231 (2023), 205-344 [Pst](https://arxiv.org/pdf/1803.01804v3): Proposition 4.36; Theorem 4.37; Proposition 4.40; pages arXiv PDF pp. 47-49; published pp. 621-624; arXiv:1803.01804v3, 10 November 2022; numbering also checked in Invent. Math. 232 (2023)
% Blueprint note: Clause 20: Write the BHS coordinates as (s,k,omega) and the project coordinates as (s,t,w). The translation is k=t-s and omega=w-(t-s); tau is renamed lambda. Hence s-omega=t-w, and the BHS differential degree (r,-1,1) becomes (r,r-1,0). This verifies the mathematical source and its stated regrading, not the outstanding Lean comparison with the fixed project model. Existing Lean proof and transitive-dependency status must remain unchanged. Spectral-sequence page and E-infinity subquotient statements here do not assert convergence to untruncated homotopy groups; BHS Convention A.5 allows an arbitrary spectrum for Theorem A.8 and Corollaries A.9/A.11. Primary references: [BHS](https://arxiv.org/pdf/1910.14116v3): Theorem A.8; Notation A.6; pages arXiv PDF pp. 71-72; published pp. 306-307; arXiv:1910.14116v3, 10 May 2022; checked against Acta Math. 231 (2023), 205-344 [BHS](https://arxiv.org/pdf/1910.14116v3): paragraph immediately before Notation A.7; proof of Theorem A.8; pages arXiv PDF p. 72; published p. 307; arXiv:1910.14116v3, 10 May 2022; checked against Acta Math. 231 (2023), 205-344
% Blueprint note: Clause 21: BHS assumes E-nilpotent completeness and a strongly convergent E-based Adams spectral sequence (Definitions 9.16–9.17). BHSObjectApplicability is the project's explicit Adams-tower encoding of these hypotheses, not a quoted definition from BHS. The mathematical source has been checked; identification with the project's specified coefficient object, Adams towers, quotient maps and realization remains part of the existing formal comparison obligations. No Lean proof obligation is discharged by this citation. Primary references: [BHS](https://arxiv.org/pdf/1910.14116v3): Theorem A.1, permanent-cycle paragraph preceding (2a); proof of Theorem 9.19(2)–(4); pages PDF pp. 70, 76–77; arXiv:1910.14116v3, 10 May 2022; final prepublication version, not published pagination [BHS](https://arxiv.org/pdf/1910.14116v3): Corollary A.19; pages PDF p. 76; arXiv:1910.14116v3, 10 May 2022; final prepublication version, not published pagination
% Blueprint note: Clause 22: Specialization to mod-two homology and the project’s precise nu and coefficient functors; no finiteness, connectivity, or completeness premise is introduced. Primary references: [Pst](https://arxiv.org/pdf/1803.01804v3): Lemma 4.23; pages 43; arXiv:1803.01804v3, 10 November 2022
% Blueprint note: Clause 23: BHS proves the source result for E-nilpotent complete spectra with a strongly convergent Adams spectral sequence. The selected family includes the ordinary sphere, its nu-cofiber, the detector, and shifts; membership alone does not establish BHS completeness. This entry is a pending source-to-selected-model comparison, not a direct application of BHS. The mathematical source has been checked; identification with the project's specified coefficient object, Adams towers, quotient maps and realization remains part of the existing formal comparison obligations. No Lean proof obligation is discharged by this citation. Primary references: [BHS](https://arxiv.org/pdf/1910.14116v3): Theorem A.1(2a), equivalently Theorem 9.19(2a); pages PDF pp. 70, 77 (also p. 36); arXiv:1910.14116v3, 10 May 2022; final prepublication version, not published pagination [BHS](https://arxiv.org/pdf/1910.14116v3): Theorem A.1(2b) and (3b), equivalently Theorem 9.19(2b),(3b); pages PDF pp. 70, 77 (also p. 36); arXiv:1910.14116v3, 10 May 2022; final prepublication version, not published pagination [BHS](https://arxiv.org/pdf/1910.14116v3): Theorem A.1(3a), equivalently Theorem 9.19(3a); pages PDF pp. 70, 77 (also p. 36); arXiv:1910.14116v3, 10 May 2022; final prepublication version, not published pagination
% Blueprint note: Clause 24: The cited papers concern the usual classical spectrum or its 2-primary/2-completed version. Recognition of the fixed project category, Adams tower, coefficient object, and selected Lean constants is a separate, currently pending formal comparison; this audit does not prove it. Primary references: [BR21](https://www.sas.rochester.edu/mth/sites/doug-ravenel/otherpapers/rognes-bruner-tmf-book.pdf): BR21 Proposition1.4;Table1.1,p56;Remark1.5,pp56–57;Figure1.10,p57; IWX Table4,p215,row(45,9,24) [IWX](https://www.numdam.org/item/10.1007/s10240-023-00139-1.pdf): BR21 Proposition1.4;Table1.1,p56;Remark1.5,pp56–57;Figure1.10,p57; IWX Table4,p215,row(45,9,24)
% Blueprint note: Clause 25: Project normalized eta factors nu eta by one lambda; realization inverts lambda and sends nu eta to the same classical eta. This verifies the mathematical source and the stated adaptation, not the pending comparison with the fixed project model. No Lean proof obligation is discharged by this audit. Primary references: [LWX126](https://arxiv.org/abs/2412.10879): Proposition 3.16; Notation 3.19; Notation 7.2; pages [18, 19, 42] [Pst](https://arxiv.org/pdf/1803.01804v3): Proposition 4.40 and Remark 4.41; pages 49
% Blueprint note: Clause 26: Write the BHS coordinates as (s,k,omega) and the project coordinates as (s,t,w). The translation is k=t-s and omega=w-(t-s); tau is renamed lambda. Hence s-omega=t-w, and the BHS differential degree (r,-1,1) becomes (r,r-1,0). This verifies the mathematical source and its stated regrading, not the outstanding Lean comparison with the fixed project model. Existing Lean proof and transitive-dependency status must remain unchanged. Spectral-sequence page and E-infinity subquotient statements here do not assert convergence to untruncated homotopy groups; BHS Convention A.5 allows an arbitrary spectrum for Theorem A.8 and Corollaries A.9/A.11. Primary references: [BHS](https://arxiv.org/pdf/1910.14116v3): Theorem A.8; Notation A.6; pages arXiv PDF pp. 71-72; published pp. 306-307; arXiv:1910.14116v3, 10 May 2022; checked against Acta Math. 231 (2023), 205-344
% Blueprint note: Clause 27: The interface existentially chooses a connecting arrow and its distinguished triangle under positive filtration and homology exactness; its full-lift comparison is only modulo lambda torsion, as in AimPaper Remark 3.18. Primary references: [BHS](https://arxiv.org/pdf/1910.14116v3): Lemma 9.15 and its proof; pages 34–35; arXiv:1910.14116v3, 10 May 2022 [LWX126](https://arxiv.org/abs/2412.10879): Propositions 3.16–3.17; Remark 3.18; pages 18–19; Repository-frozen AimPaper/main.tex and AimPaper/2412.10879.pdf
% Blueprint note: Clause 28: The exceptional Hopf map is explicitly the degree-two map. This verifies the mathematical source and the stated adaptation, not the pending comparison with the fixed project model. No Lean proof obligation is discharged by this audit. Primary references: [Adams58](https://www.sas.rochester.edu/mth/sites/doug-ravenel/otherpapers/Adams-SS.pdf): Lemma 2.3 and the following m=0 paragraph; pages 184
% Blueprint note: Clause 29: The cited papers concern the usual classical spectrum or its 2-primary/2-completed version. Recognition of the fixed project category, Adams tower, coefficient object, and selected Lean constants is a separate, currently pending formal comparison; this audit does not prove it. Primary references: [BR21](https://www.sas.rochester.edu/mth/sites/doug-ravenel/otherpapers/rognes-bruner-tmf-book.pdf): Proposition1.4;Table1.1,p54;Figure1.9,p56
% Blueprint note: Clause 30: Equal h5^2 detection implies difference in F3; absence of graded pieces3,4,5 gives F6. This verifies the mathematical source and the stated adaptation, not the pending comparison with the fixed project model. No Lean proof obligation is discharged by this audit. Primary references: [LWX126](https://arxiv.org/abs/2412.10879): Proof of Lemma 7.10; pages 45 [IWXCharts](https://zenodo.org/records/6987157/files/Adams-classical-Einfty.csv?download=1): Adams-classical-Einfty.csv, physical rows 204--207 (including header)
% Blueprint note: Clause 31: BHS assumes E-nilpotent completeness and a strongly convergent E-based Adams spectral sequence (Definitions 9.16–9.17). BHSObjectApplicability is the project's explicit Adams-tower encoding of these hypotheses, not a quoted definition from BHS. The mathematical source has been checked; identification with the project's specified coefficient object, Adams towers, quotient maps and realization remains part of the existing formal comparison obligations. No Lean proof obligation is discharged by this citation. Primary references: [BHS](https://arxiv.org/pdf/1910.14116v3): Theorem A.1, permanent-cycle paragraph preceding (2a); proof of Theorem 9.19(2)–(4); pages PDF pp. 70, 76–77; arXiv:1910.14116v3, 10 May 2022; final prepublication version, not published pagination [BHS](https://arxiv.org/pdf/1910.14116v3): Corollary A.19; pages PDF p. 76; arXiv:1910.14116v3, 10 May 2022; final prepublication version, not published pagination
% Blueprint note: Clause 32: The cited papers concern the usual classical spectrum or its 2-primary/2-completed version. Recognition of the fixed project category, Adams tower, coefficient object, and selected Lean constants is a separate, currently pending formal comparison; this audit does not prove it. Primary references: [BR21](https://www.sas.rochester.edu/mth/sites/doug-ravenel/otherpapers/rognes-bruner-tmf-book.pdf): Section0.1,p1;Section0.5,pp8–9;Section1.2,p64
% Blueprint note: Clause 33: The current type asserts a distinguished triangle on unshifted nu(X), nu(Y), nu(Z). Proposition 3.20 retains nontrivial weight shifts. The source does not directly justify removing them; no fixed-model counterexample or interface change is asserted here. Primary references: [LWX126](https://arxiv.org/abs/2412.10879): Notation 3.19; Proposition 3.20; pages 19; Repository-frozen AimPaper/main.tex and AimPaper/2412.10879.pdf
% Blueprint note: Clause 34: The exact interface is the signed TC3 pushpull consequence for these selected tensor/suspension conventions, retaining the common vertex, all comparison squares, and represented homotopy lift. Primary references: [May01](https://www.math.uchicago.edu/~may/PAPERS/AddJan01.pdf): Axiom TC3; Lemma 4.6; pages 12–14; Author-hosted January 2001 preprint, 31 pages; numbered statements retain published numbering
% Blueprint note: Clause 35: Write the BHS coordinates as (s,k,omega) and the project coordinates as (s,t,w). The translation is k=t-s and omega=w-(t-s); tau is renamed lambda. Hence s-omega=t-w, and the BHS differential degree (r,-1,1) becomes (r,r-1,0). This verifies the mathematical source and its stated regrading, not the outstanding Lean comparison with the fixed project model. Existing Lean proof and transitive-dependency status must remain unchanged. Spectral-sequence page and E-infinity subquotient statements here do not assert convergence to untruncated homotopy groups; BHS Convention A.5 allows an arbitrary spectrum for Theorem A.8 and Corollaries A.9/A.11. Primary references: [Pst](https://arxiv.org/pdf/1803.01804v3): Proposition 4.36; Theorem 4.37; Proposition 4.40; pages arXiv PDF pp. 47-49; published pp. 621-624; arXiv:1803.01804v3, 10 November 2022; numbering also checked in Invent. Math. 232 (2023) [BHS](https://arxiv.org/pdf/1910.14116v3): paragraph immediately before Notation A.7; proof of Theorem A.8; pages arXiv PDF p. 72; published p. 307; arXiv:1910.14116v3, 10 May 2022; checked against Acta Math. 231 (2023), 205-344 [BHS](https://arxiv.org/pdf/1910.14116v3): Theorem A.8; Notation A.6; pages arXiv PDF pp. 71-72; published pp. 306-307; arXiv:1910.14116v3, 10 May 2022; checked against Acta Math. 231 (2023), 205-344
% Blueprint note: Clause 36: The cited papers concern the usual classical spectrum or its 2-primary/2-completed version. Recognition of the fixed project category, Adams tower, coefficient object, and selected Lean constants is a separate, currently pending formal comparison; this audit does not prove it. Primary references: [BR21](https://www.sas.rochester.edu/mth/sites/doug-ravenel/otherpapers/rognes-bruner-tmf-book.pdf): BR21 Proposition1.4;Table1.1,p56;Figure1.10,p57; IWX Table4,p215,row(45,9,24) [IWX](https://www.numdam.org/item/10.1007/s10240-023-00139-1.pdf): BR21 Proposition1.4;Table1.1,p56;Figure1.10,p57; IWX Table4,p215,row(45,9,24)
% Blueprint note: Detailed source attribution: (Adapted from the individually cited results below. The selected-model comparisons, the extension to $N<2$, and the unshifted conclusion retain the stated verification limits; see \cite{Pst,BHS,LWX126,BR21,Moss,May01}.)
% Blueprint note: Clause 1 detailed source attribution: (Adapted from \cite{Xu}, Corollary 1.3, p.~1612; and \cite{IWX}, Corollary 1.3 and Table 1, row 62.)
% Blueprint note: Clause 2 detailed source attribution: (Adapted from \cite{Pst}, Proposition 4.40, and \cite{BHS}, Theorem A.8(2) and the comparison before Notation A.7.)
% Blueprint note: Clause 3 detailed source attribution: (Adapted from \cite{SchwedeSymmetricSpectra}, Proposition II.5.22 and Definition II.6.21, specialized to $H\mathbb F_2$; cf.\ \cite{EKMM}, Proposition IV.1.2(i).)
% Blueprint note: Clause 4 detailed source attribution: (Adapted from \cite{Adams58}, Lemma 2.3, p.~184; and \cite{Adams60}, Section 1.1, p.~20.)
% Blueprint note: Clause 5 detailed source attribution: (Adapted from \cite{BHS}, Theorem A.8(2).)
% Blueprint note: Clause 6 detailed source attribution: (Adapted from \cite{BHS}, Theorem A.1(2a), equivalently Theorem 9.19(2a).)
% Blueprint note: Clause 7 detailed source attribution: (Adapted from \cite{BHS}, Theorem A.1(2a),(2b),(3a),(3b). Source-to-selected-model comparison pending.)
% Blueprint note: Clause 8 detailed source attribution: (Adapted from \cite{BHS}, Lemma 9.15, as restated in \cite{LWX126}, Proposition 3.16; comparison with the selected Adams tower remains pending.)
% Blueprint note: Clause 9 detailed source attribution: (Adapted from \cite{BHS}, Theorem A.1(2b),(3b), equivalently Theorem 9.19(2b),(3b).)
% Blueprint note: Clause 10 detailed source attribution: (Adapted from \cite{BHS}, Theorem A.1(3a), equivalently Theorem 9.19(3a).)
% Blueprint note: Clause 11 detailed source attribution: (Adapted from \cite{BHS}, Theorem A.8 and its proof.)
% Blueprint note: Clause 12 detailed source attribution: (Adapted from \cite{BHS}, Theorem A.8 and the comparison before Notation A.7.)
% Blueprint note: Clause 13 detailed source attribution: (Adapted from \cite{BR21}, Proposition 1.4 and Table 1.1, p. 54, and Remark 1.5, pp. 56--57.)
% Blueprint note: Clause 14 detailed source attribution: (Adapted from \cite{BHS}, Theorem A.1(2b),(3b), equivalently Theorem 9.19(2b),(3b).)
% Blueprint note: Clause 15 detailed source attribution: (Adapted from \cite{Moss}, Section 1, pp. 283--285; the tower formulation follows \cite{BelmontKong2024}, Definition 2.3.)
% Blueprint note: Clause 16 detailed source attribution: (Adapted from \cite{Adams58}, Lemma 2.3, p.~184; and \cite{Adams60}, Section 1.1, p.~20.)
% Blueprint note: Clause 17 detailed source attribution: (Adapted from \cite{BR21}, Section 0.3, p. 5, Section 0.5, pp. 8--9, and Section 1.2, p. 64.)
% Blueprint note: Clause 18 detailed source attribution: (Adapted from \cite{SchwedeSymmetricSpectra}, Example I.1.8, p.~12, and the proof of Proposition II.5.21, p.~264.)
% Blueprint note: Clause 19 detailed source attribution: (Adapted from \cite{BHS}, the comparison before Notation A.7 and the proof of Theorem A.8; \cite{Pst}, Proposition 4.40.)
% Blueprint note: Clause 20 detailed source attribution: (Adapted from \cite{BHS}, Theorem A.8 and the comparison before Notation A.7.)
% Blueprint note: Clause 21 detailed source attribution: (Adapted from \cite{BHS}, Theorem A.1, the permanent-cycle paragraph, and Corollary A.19.)
% Blueprint note: Clause 22 detailed source attribution: (Adapted from \cite{Pst}, Lemma 4.23, as restated in \cite{LWX126}, Proposition 3.1; recognition of the selected synthetic model remains pending.)
% Blueprint note: Clause 23 detailed source attribution: (Adapted from \cite{BHS}, Theorem A.1(2a),(2b),(3a),(3b). Source-to-selected-model comparison pending.)
% Blueprint note: Clause 24 detailed source attribution: (Adapted from \cite{BR21}, Proposition 1.4 and Table 1.1, p. 56, and Remark 1.5, pp. 56--57; notation as in \cite{IWX}, Table 4, p. 215.)
% Blueprint note: Clause 25 detailed source attribution: (Adapted from \cite{LWX126}, Proposition 3.16 and Notation 3.19; and \cite{Pst}, Proposition 4.40 and Remark 4.41.)
% Blueprint note: Clause 26 detailed source attribution: (Adapted from \cite{BHS}, Theorem A.8(3). The extension to $N<2$ remains unverified.)
% Blueprint note: Clause 27 detailed source attribution: (Adapted from \cite{LWX126}, Proposition 3.17 and Remark 3.18, using the proof of \cite{BHS}, Lemma 9.15; comparison with the selected connecting map remains pending.)
% Blueprint note: Clause 28 detailed source attribution: (Adapted from \cite{Adams58}, Lemma 2.3 and the following m=0 paragraph, p.~184.)
% Blueprint note: Clause 29 detailed source attribution: (Adapted from \cite{BR21}, Proposition 1.4 and Table 1.1, p. 54, and Figure 1.9, p. 56.)
% Blueprint note: Clause 30 detailed source attribution: (From \cite{LWX126}, proof of Lemma 7.10, p.~45; independently checked against \cite{IWXCharts}, the stem-62 E-infinity rows.)
% Blueprint note: Clause 31 detailed source attribution: (Adapted from \cite{BHS}, Theorem A.1, the permanent-cycle paragraph, and Corollary A.19.)
% Blueprint note: Clause 32 detailed source attribution: (Adapted from \cite{BR21}, Section 0.1, p. 1, and Section 1.2, p. 64.)
% Blueprint note: Clause 33 detailed source attribution: (Related source: \cite{LWX126}, Proposition 3.20, which retains weight shifts; the present unshifted conclusion is not verified.)
% Blueprint note: Clause 34 detailed source attribution: (Adapted from \cite{May01}, Axiom TC3 and Lemma 4.6; see \cite{LWX126}, Lemma 6.11. The signed comparison for the selected tensor context remains pending.)
% Blueprint note: Clause 35 detailed source attribution: (Adapted from \cite{Pst}, Proposition 4.40, and \cite{BHS}, the comparison before Notation A.7 and Theorem A.8.)
% Blueprint note: Clause 36 detailed source attribution: (Adapted from \cite{BR21}, Proposition 1.4 and Table 1.1, p. 56, and Figure 1.10, p. 57; notation as in \cite{IWX}, Table 4, p. 215.)
\end{theorem}

\begin{theorem}[The Bockstein represents the classical differential]
\label{frontier:Literature:98ed1948c077}
\level{fine}
(Adapted from \cite{BHS}, Thm. A.1(1c).)\par{}

Let $X$ belong to the chosen classical family, $q>0$, and $x\in E_2^{s,t}(X)$ reach $E_{q+1}$. There exists a lift $a\in\pi_{t-s,t}(\nu X/\lambda^q)$ restricting to the first-quotient class of $x$ such that the Bockstein label of $a$ represents $d_{q+1}x\in E_{q+1}^{s+q+1,t+q}(X)$.
\lean{KIP126.Literature.Standard.bhs_bockstein}
% Blueprint note: Source verification (2026-10-18): Source adaptation checked. Finite-quotient result: write q for BHS r, identify tau with lambda, and use stem k=t-s. Reaching E_(q+1) means d_2 through d_q vanish; zero and earlier boundaries are allowed. Unlike the untruncated conclusion, the finite case extends to arbitrary spectra via Corollary A.11 and Lemma A.14, as in the proof of Theorem A.1(1). Over F_2 the sign in A.1(1c) disappears. The mathematical source has been checked; identification with the project's specified coefficient object, Adams towers, quotient maps and realization remains part of the existing formal comparison obligations. No Lean proof obligation is discharged by this citation. Primary references: [BHS](https://arxiv.org/pdf/1910.14116v3): Theorem A.1(1a)–(1c); Corollary A.11; Lemma A.14; proof of Theorem 9.19(1); pages PDF pp. 70, 73–75; arXiv:1910.14116v3, 10 May 2022; final prepublication version, not published pagination [LWX126](https://arxiv.org/abs/2412.10879): Theorem 3.8(1)–(2); pages PDF p. 16; Repository AimPaper/main.tex and AimPaper/2412.10879.pdf; exact local source snapshot
% Blueprint note: Detailed source attribution: (Adapted from \cite{BHS}, Theorem A.1(1c) and its proof, Corollary A.11 and Lemma A.14; \cite{LWX126}, Theorem 3.8(2).)
\end{theorem}



\begin{theorem}[A synthetic lift with prescribed realization]
\label{frontier:Literature:acb86a0972cf}
\level{fine}
(Adapted from \cite{BHS}, Thm. A.1(3b).)\par{}

For $X$ in the chosen classical family, let $x\in E_2^{s,t}(X)$ be a nonzero permanent label detecting $a\in\pi_{t-s}X$. There exists $b\in\pi_{t-s,t}\nu X$ with first-quotient label $x$ and realization equal to $a$.
\lean{KIP126.Literature.Standard.selected_realization_prescribed_lift}
% Blueprint note: Source verification (2026-10-18): Partial source correspondence. BHS proves the source result for E-nilpotent complete spectra with a strongly convergent Adams spectral sequence. The selected family includes the ordinary sphere, its nu-cofiber, the detector, and shifts; membership alone does not establish BHS completeness. This entry is a pending source-to-selected-model comparison, not a direct application of BHS. The mathematical source has been checked; identification with the project's specified coefficient object, Adams towers, quotient maps and realization remains part of the existing formal comparison obligations. No Lean proof obligation is discharged by this citation. Primary references: [BHS](https://arxiv.org/pdf/1910.14116v3): Theorem A.1(2b) and (3b), equivalently Theorem 9.19(2b),(3b); pages PDF pp. 70, 77 (also p. 36); arXiv:1910.14116v3, 10 May 2022; final prepublication version, not published pagination
% Blueprint note: Detailed source attribution: (Adapted from \cite{BHS}, Theorem A.1(3b). Source-to-selected-model comparison pending.)
\end{theorem}

\begin{theorem}[Projection and lambda on infinite-page windows]
\label{frontier:Literature:b8fc8580677d}
\level{fine}
(Adapted from \cite{BHS}, Cors. A.9 and A.11.)\par{}

Let $X$ belong to the chosen classical family.
\begin{enumerate}
\item For $q>0$ and $0\le t-w<q$, projection from $\nu X$ to $\nu X/\lambda^q$ is the natural map from $Z_\infty^{s,t}/B_{1+t-w}^{s,t}$ to $Z_{q-t+w}^{s,t}/B_{1+t-w}^{s,t}$.
\item For $k\ge0$ and $w\le t$, multiplication by $\lambda^k$ on $E_\infty(\nu X)$ is the quotient map enlarging $B_{1+t-w}^{s,t}$ to $B_{1+t-w+k}^{s,t}$.
\item For $0\le k<q$ and $0\le t-w<q-k$, the map $\lambda^k$ from length $q-k$ and weight $w$ to length $q$ and weight $w-k$ is the natural cycle/boundary quotient map under the finite windows.
\end{enumerate}
\lean{KIP126.Literature.Standard.eInfinity_rho_nu,KIP126.Literature.Standard.eInfinity_lambda_nu,KIP126.Literature.Standard.eInfinity_lambda_finite}
% Blueprint note: Source verification (2026-10-18): Source adaptation checked. BHS Corollary A.9 gives the untruncated E-infinity subquotient, and Corollary A.11 gives the finite quotient. All formula indices transform by s-omega=t-w. The finite formula holds for q>=1 in BHS, while AimPaper Proposition 3.12 explicitly states q>=2. Additive equivalences are viewed as Z-linear without a change of scalar structure. Matching a specific E2 representative to its quotient class requires the canonical label-compatible presentation, not an arbitrary abstract isomorphism. Zero labels are permitted in subgroup and quotient statements. Write the BHS coordinates as (s,k,omega) and the project coordinates as (s,t,w). The translation is k=t-s and omega=w-(t-s); tau is renamed lambda. Hence s-omega=t-w, and the BHS differential degree (r,-1,1) becomes (r,r-1,0). This verifies the mathematical source and its stated regrading, not the outstanding Lean comparison with the fixed project model. Existing Lean proof and transitive-dependency status must remain unchanged. Spectral-sequence page and E-infinity subquotient statements here do not assert convergence to untruncated homotopy groups; BHS Convention A.5 allows an arbitrary spectrum for Theorem A.8 and Corollaries A.9/A.11. Primary references: [BHS](https://arxiv.org/pdf/1910.14116v3): Corollary A.9(1)-(4), including its final surjectivity sentence; pages arXiv PDF pp. 72-73; published pp. 307-308; arXiv:1910.14116v3, 10 May 2022; checked against Acta Math. 231 (2023), 205-344 [BHS](https://arxiv.org/pdf/1910.14116v3): Corollary A.11(1)-(3) and its proof; Notation A.10; pages arXiv PDF p. 73; published p. 308; arXiv:1910.14116v3, 10 May 2022; checked against Acta Math. 231 (2023), 205-344 [LWX126](https://arxiv.org/html/2412.10879v2#S3): Remark 3.14; pages local PDF p. 17; arXiv:2412.10879v2, 22 February 2025; local AimPaper/main.tex and PDF also checked
% Blueprint note: Detailed source attribution: (Adapted from \cite{BHS}, Corollaries A.9 and A.11, and \cite{LWX126}, Remark 3.14.)
\end{theorem}



\begin{theorem}[The order and label of BX theta-five]
\label{frontier:Literature:cb4dd3fedc28}
\level{fine}
(Adapted from \cite{BurklundXu}, Prop. 7.19.)\par{}

The distinguished class $\theta_5^{\mathrm{BX}}\in\pi_{62,64}S^{0,0}$ satisfies $2\theta_5^{\mathrm{BX}}=0$ and has first-quotient label $h_5^2$.
\lean{KIP126.Literature.Standard.bx_order_two,KIP126.Literature.Standard.bx_theta_label}
% Blueprint note: Source verification (2026-10-18): Source adaptation checked. This is the project-selected BX source choice, not a claim for every class with the same label. BX uses bidegree (stem, Adams filtration); the project uses (stem, weight), with weight = stem + filtration. Thus the source theta in degree (62,2) is displayed here in degree (62,64), and the source total boundary target is regraded to the project shift (1,-1). The finite statement retains r >= 1. Source-to-model identification and the registered Lean proof obligations remain separate from this citation check. Primary references: [BurklundXu](https://arxiv.org/pdf/2302.11869v3): Proposition 7.19 and its proof; Theorem 7.20 and its proof; pages 54; arXiv:2302.11869v3, 9 November 2024; accepted version [LWX126](https://arxiv.org/abs/2412.10879): Theorem 7.3, Remarks 7.4–7.5; Example 5.8(4); proof of Lemma 7.16; pages 42, 28, 49–50; Repository AimPaper/main.tex and AimPaper/2412.10879.pdf snapshot; no unstated version substitution
% Blueprint note: Detailed source attribution: (Adapted from \cite{BurklundXu}, Proposition 7.19 and the proof of Theorem 7.20; see \cite{LWX126}, Notation 7.2 and Remark 7.5.)
\end{theorem}

\begin{theorem}[Detection of the classical Hopf eta]
\label{frontier:Literature:da6f8834fbe5}
\level{fine}
(Adapted from \cite{Adams58}, Lem. 2.3.)\par{}

The label $h_1\in E_2^{1,2}(S^0)$ detects the classical Hopf map $\eta\in\pi_1S^0$.
\lean{KIP126.Literature.Standard.eta_detection}
% Blueprint note: Source verification (2026-10-18): Source adaptation checked. This fixes the actual classical Hopf eta and the standard h1 label on the project Adams sequence. This verifies the mathematical source and the stated adaptation, not the pending comparison with the fixed project model. No Lean proof obligation is discharged by this audit. Primary references: [Adams58](https://www.sas.rochester.edu/mth/sites/doug-ravenel/otherpapers/Adams-SS.pdf): Lemma 2.3 and the following m=0 paragraph; pages 184 [Adams60](https://www.maths.ed.ac.uk/~v1ranick/papers/adams1.pdf): Section 1.1, after Theorem 1.1.1; pages 20
% Blueprint note: Detailed source attribution: (Adapted from \cite{Adams58}, Lemma 2.3, p.~184; and \cite{Adams60}, Section 1.1, p.~20.)
\end{theorem}

\begin{theorem}[Infinite representatives in a finite window]
\label{frontier:Literature:dd88947b5ffc}
\level{fine}
(Adapted from \cite{BHS}, Cor. A.11.)\par{}

Let $X$ belong to the chosen classical family, $q>0$, $0\le t-w<q$, $x\in E_2^{s,t}(X)$ and $z\in E_\infty^{s,t,w}(\nu X/\lambda^q)$. The projected weight-$w$ label of $x$ represents $z$ if and only if $x\in Z_{q-t+w}^{s,t}(X)$ and the finite-window image of $z$ is the class of $x$ modulo $B_{1+t-w}^{s,t}(X)$.
\lean{KIP126.Literature.Standard.eInfinity_finite_labels}
% Blueprint note: Source verification (2026-10-18): Source adaptation checked. BHS Corollary A.9 gives the untruncated E-infinity subquotient, and Corollary A.11 gives the finite quotient. All formula indices transform by s-omega=t-w. The finite formula holds for q>=1 in BHS, while AimPaper Proposition 3.12 explicitly states q>=2. Additive equivalences are viewed as Z-linear without a change of scalar structure. Matching a specific E2 representative to its quotient class requires the canonical label-compatible presentation, not an arbitrary abstract isomorphism. Zero labels are permitted in subgroup and quotient statements. Write the BHS coordinates as (s,k,omega) and the project coordinates as (s,t,w). The translation is k=t-s and omega=w-(t-s); tau is renamed lambda. Hence s-omega=t-w, and the BHS differential degree (r,-1,1) becomes (r,r-1,0). This verifies the mathematical source and its stated regrading, not the outstanding Lean comparison with the fixed project model. Existing Lean proof and transitive-dependency status must remain unchanged. Spectral-sequence page and E-infinity subquotient statements here do not assert convergence to untruncated homotopy groups; BHS Convention A.5 allows an arbitrary spectrum for Theorem A.8 and Corollaries A.9/A.11. Primary references: [BHS](https://arxiv.org/pdf/1910.14116v3): Corollary A.11(1)-(3) and its proof; Notation A.10; pages arXiv PDF p. 73; published p. 308; arXiv:1910.14116v3, 10 May 2022; checked against Acta Math. 231 (2023), 205-344 [LWX126](https://arxiv.org/html/2412.10879v2#S3): Proposition 3.12; pages local PDF p. 17; arXiv:2412.10879v2, 22 February 2025; local AimPaper/main.tex and PDF also checked [LWX126](https://arxiv.org/html/2412.10879v2#S3): Remark 3.14; pages local PDF p. 17; arXiv:2412.10879v2, 22 February 2025; local AimPaper/main.tex and PDF also checked
% Blueprint note: Detailed source attribution: (Adapted from \cite{BHS}, Corollary A.11, and \cite{LWX126}, Proposition 3.12 and Remark 3.14.)
\end{theorem}

\begin{theorem}[Permanent representatives and untruncated lifts]
\label{frontier:Literature:e0657653596f}
\level{fine}
(Adapted from \cite{BHS}, Thm. A.1.)\par{}

For $X$ in the chosen classical family and $x\in E_2^{s,t}(X)$, the label $x$ has a common infinite-cycle representative if and only if there exists $a\in\pi_{t-s,t}\nu X$ with first-quotient label $x$. Zero labels are allowed.
\lean{KIP126.Literature.Standard.selected_permanent_lift}
% Blueprint note: Source verification (2026-10-18): Partial source correspondence. BHS proves the source result for E-nilpotent complete spectra with a strongly convergent Adams spectral sequence. The selected family includes the ordinary sphere, its nu-cofiber, the detector, and shifts; membership alone does not establish BHS completeness. This entry is a pending source-to-selected-model comparison, not a direct application of BHS. The mathematical source has been checked; identification with the project's specified coefficient object, Adams towers, quotient maps and realization remains part of the existing formal comparison obligations. No Lean proof obligation is discharged by this citation. Primary references: [BHS](https://arxiv.org/pdf/1910.14116v3): Theorem A.1, permanent-cycle paragraph preceding (2a); proof of Theorem 9.19(2)–(4); pages PDF pp. 70, 76–77; arXiv:1910.14116v3, 10 May 2022; final prepublication version, not published pagination
% Blueprint note: Detailed source attribution: (Adapted from \cite{BHS}, Theorem A.1, the permanent-cycle paragraph and its proof on pp. 76--77 (arXiv v3). Source-to-selected-model comparison pending.)
\end{theorem}

\begin{theorem}[Quotient algebra laws and the first-quotient product]
\label{frontier:Literature:a0113c9a3b80}
\level{fine}
(Adapted from \cite{Pst}, Thm. 4.46.)\par{}

Let $s,t,s^{\prime},t^{\prime}\in\mathbb N$ and let $x,y$ be first-quotient classes in bidegrees $(t-s,t)$ and $(t^{\prime}-s^{\prime},t^{\prime})$. The first-quotient image of their algebra product equals the Ext product of their images, in bidegree $(s+s^{\prime},t+t^{\prime})$.

(Adapted from \cite{LWX126}, Remark 3.4.)\par{}

For every $q>0$, the chosen algebra on $S^{0,0}/\lambda^q$ is commutative, and its unit is the quotient inclusion $S^{0,0}\to S^{0,0}/\lambda^q$.
\lean{KIP126.Literature.Standard.quotientAlgebra_commutative,KIP126.Literature.Standard.quotientAlgebra_first_quotient,KIP126.Literature.Standard.quotientAlgebra_unit}
% Blueprint note: Source verification (2026-10-18): Source adaptation checked. Symmetric monoidality gives multiplicativity of the canonical Ext comparison, with s,t,s-prime,t-prime nonnegative and stem t-s. Identification with the project’s selected firstQuotient and sphereAlgebra product is still an explicit pending comparison. Primary references: [Pst](https://arxiv.org/pdf/1803.01804v3): Corollary 4.30; Theorem 4.46; Lemma 4.56; pages 46, 50–51, 54; arXiv:1803.01804v3, 10 November 2022
% Blueprint note: Detailed source attribution: (Adapted from \cite{Pst}, Theorem 4.46 and Lemma 4.56.)
% Blueprint note: Source verification (2026-10-18): Source adaptation checked. A source E-infinity algebra yields a commutative monoid in the homotopy category with the quotient map as unit. Matching the selected MonObj and its actual unit to that source structure remains a comparison obligation. Primary references: [LWX126](https://arxiv.org/abs/2412.10879): Remark 3.4; pages 15; Repository-frozen AimPaper/main.tex and AimPaper/2412.10879.pdf [BHSmot](https://arxiv.org/pdf/2010.10325v2): Definition B.3; Lemma B.5; Proposition C.2; Proposition C.22; pages 69–70, 75; arXiv:2010.10325v2, 30 June 2022 [Pst](https://arxiv.org/pdf/1803.01804v3): Corollary 4.30; Theorem 4.46; Lemma 4.56; pages 46, 50–51, 54; arXiv:1803.01804v3, 10 November 2022
% Blueprint note: Detailed source attribution: (Adapted from \cite{LWX126}, Remark 3.4; for the first quotient see \cite{Pst}, Corollary 4.30.)
\end{theorem}

\begin{theorem}[Low Toda comparisons and the two-primary stem]
\label{frontier:Literature:dd011cab6d38}
\level{fine}
(Adapted from \cite{LWX126}, proof of Lem. 7.16.)\par{}

In the synthetic sphere, $\eta^2\in\langle[h_0],\eta,[h_0]\rangle$. Every $a\in\pi_{2,3}S^{0,0}$ satisfies $[h_0]a=0$.

(Adapted from \cite{Toda1962}, Cor. 3.7; \cite{LWX126}, proof of Lem. 7.16.)\par{}

For every $\theta\in\pi_{62,64}S^{0,0}$ with $2\theta=0$, one has $\lambda^2\eta\theta\in\langle2,\theta,2\rangle$. The first-quotient label of $[h_0]$ is $h_0$. If $a\in\pi_{62}S^0$ is annihilated by a power of two, then $2a=0$.
\lean{KIP126.Literature.Standard.stem62_exponent_two,KIP126.Literature.Standard.toda_eta_squared,KIP126.Literature.Standard.toda_h0_label,KIP126.Literature.Standard.toda_low_indeterminacy,KIP126.Literature.Standard.toda_symmetric}
% Blueprint note: Source verification (2026-10-18): Source adaptation checked. The fixed h-zero choice is normalized by the first-quotient label and lambda[h-zero]=2, as in AimPaper Example 5.8(4). BHS Proposition A.20 gives a complete low synthetic ring presentation, with relation (0) h-zero eta=0 and relation (9) tau h-zero=2. In stem two, the presentation forces pi_(2,3) to be generated by tau eta^2 (the h-zero factors vanish), so h-zero annihilates every element there. This last assertion is a derived low-degree consequence, not a verbatim statement of AimPaper Lemma 7.16. The project must still identify its selected constants with these normalized low generators. Primary references: [LWX126](https://arxiv.org/abs/2412.10879): Theorem 7.3, Remarks 7.4–7.5; Example 5.8(4); proof of Lemma 7.16; pages 42, 28, 49–50; Repository AimPaper/main.tex and AimPaper/2412.10879.pdf snapshot; no unstated version substitution [BHS](https://arxiv.org/pdf/1910.14116v3): Remark 9.5; Theorem A.8; Corollary A.9(1); Notation A.10, Corollary A.11; Proposition A.20, relations (0), (9), and Remark A.22; pages 32, 72–73, 78–79; published Acta page mapping supplied in bhs_pages/proposals.json; arXiv:1910.14116v3, as independently retrieved by the BHS audit worker
% Blueprint note: Detailed source attribution: (Adapted from \cite{LWX126}, proof of Lemma 7.16, and \cite{BHS}, Proposition A.20.)
% Blueprint note: Source verification (2026-10-18): Partial source correspondence. Toda Corollary 3.7 gives the uniform classical order-two symmetric-bracket law. Baues--Muro Definition 1.2, relation (T6), and Theorem 1.3 supply the connective-ring-spectrum version with the image of the stable Hopf element. To apply it in the synthetic category, one uses the canonical stable enrichment (or the corresponding square-zero mapping module) and identifies that Hopf element in pi_(1,0) with lambda^2 eta using BHS Proposition A.20. AimPaper proof of Lemma 7.16 uses the displayed synthetic formula for theta-five. The bound Lean statement is broader: every theta in pi_(62,64) killed by two, without a theta-five detection hypothesis. The cited uniform law explains this extension, but the canonical-enrichment and selected-Hopf-class comparison is an additional source-to-project obligation, not a quoted assertion of Lemma 7.16. Primary references: [Toda1962](https://webhomes.maths.ed.ac.uk/~v1ranick/papers/toda.pdf): Corollary 3.7; pages 31 (right-hand page of PDF page 18); Composition Methods in Homotopy Groups of Spheres, original book scan [BauesMuro2009](https://idus.us.es/server/api/core/bitstreams/57be0e00-a434-4786-b7e3-83a7eeed4e7a/content): Definition 1.2, relation (T6), and Theorem 1.3; pages 3 in the author manuscript; Author manuscript arXiv:math/0612450v1, 15 December 2006; published Comm. Algebra 37 (2009), 56–82 [LWX126](https://arxiv.org/abs/2412.10879): Theorem 7.3, Remarks 7.4–7.5; Example 5.8(4); proof of Lemma 7.16; pages 42, 28, 49–50; Repository AimPaper/main.tex and AimPaper/2412.10879.pdf snapshot; no unstated version substitution [BHS](https://arxiv.org/pdf/1910.14116v3): Remark 9.5; Theorem A.8; Corollary A.9(1); Notation A.10, Corollary A.11; Proposition A.20, relations (0), (9), and Remark A.22; pages 32, 72–73, 78–79; published Acta page mapping supplied in bhs_pages/proposals.json; arXiv:1910.14116v3, as independently retrieved by the BHS audit worker [IWX](https://arxiv.org/pdf/2001.04511v3): Corollary 1.3, Table 1, row k=62; pages [4, 6]
% Blueprint note: Detailed source attribution: (Adapted from \cite{Toda1962}, Corollary 3.7; \cite{LWX126}, Example 5.8(4) and proof of Lemma 7.16; and \cite{IWX}, Corollary 1.3 and Table 1, row 62. The uniform synthetic comparison remains to be checked.)
\end{theorem}

\begin{theorem}[Hopf relations and the realization kernel]
\label{frontier:Literature:e10e68418839}
\level{fine}
(Adapted from \cite{BHS}, Prop. A.20(9).)\par{}

One has $\lambda[h_0]=2$ in the synthetic sphere. For every object $X$ in the chosen synthetic family and $a\in\pi_{m,w}X$, realization sends $a$ to zero if and only if $\lambda^k a=0$ for some $k\in\mathbb N$.

(Adapted from \cite{BHS}, Prop. A.20(0).)\par{}

In the synthetic sphere, $[h_0]\eta=0$.
\lean{KIP126.Literature.Standard.selected_realization_kernel,KIP126.Literature.Standard.toda_h0_eta,KIP126.Literature.Standard.toda_lambda_h0}
% Blueprint note: Source verification (2026-10-18): Partial source correspondence. The lambda[h_0]=2 clause is explicitly stated in AimPaper Example 5.8(4), and in the 2-complete Toda-range presentation BHS Proposition A.20(9). The realization-kernel clause follows from the full-category localization telescope of Pstragowski, not from BHS Appendix A alone. Applying that kernel statement to the project realization remains an admitted comparison; compactness of a hypercomplete sphere must not be silently assumed. The mathematical source has been checked; identification with the project's specified coefficient object, Adams towers, quotient maps and realization remains part of the existing formal comparison obligations. No Lean proof obligation is discharged by this citation. Primary references: [BHS](https://arxiv.org/pdf/1910.14116v3): Proposition A.20, relation (9); pages PDF p. 78; arXiv:1910.14116v3, 10 May 2022; final prepublication version, not published pagination [LWX126](https://arxiv.org/abs/2412.10879): Example 5.8(4); pages PDF p. 28; main.tex line 1423; Repository AimPaper/main.tex and AimPaper/2412.10879.pdf; exact local source snapshot [Pst](https://arxiv.org/pdf/1803.01804v3): Proposition 4.33; Proposition 4.40; Remark 4.41; pages PDF pp. 46–49; arXiv:1803.01804v3, 10 November 2022; not published pagination
% Blueprint note: Detailed source attribution: (Adapted from \cite{BHS}, Proposition A.20(9), and \cite{LWX126}, Example 5.8(4); \cite{Pst}, Remark 4.41. Realization-kernel model comparison pending.)
% Blueprint note: Source verification (2026-10-18): Source adaptation checked. The fixed h-zero choice is normalized by the first-quotient label and lambda[h-zero]=2, as in AimPaper Example 5.8(4). BHS Proposition A.20 gives a complete low synthetic ring presentation, with relation (0) h-zero eta=0 and relation (9) tau h-zero=2. In stem two, the presentation forces pi_(2,3) to be generated by tau eta^2 (the h-zero factors vanish), so h-zero annihilates every element there. This last assertion is a derived low-degree consequence, not a verbatim statement of AimPaper Lemma 7.16. The project must still identify its selected constants with these normalized low generators. Primary references: [LWX126](https://arxiv.org/abs/2412.10879): Theorem 7.3, Remarks 7.4–7.5; Example 5.8(4); proof of Lemma 7.16; pages 42, 28, 49–50; Repository AimPaper/main.tex and AimPaper/2412.10879.pdf snapshot; no unstated version substitution [BHS](https://arxiv.org/pdf/1910.14116v3): Remark 9.5; Theorem A.8; Corollary A.9(1); Notation A.10, Corollary A.11; Proposition A.20, relations (0), (9), and Remark A.22; pages 32, 72–73, 78–79; published Acta page mapping supplied in bhs_pages/proposals.json; arXiv:1910.14116v3, as independently retrieved by the BHS audit worker
% Blueprint note: Detailed source attribution: (Adapted from \cite{BHS}, Proposition A.20, relation (0); also implicit in \cite{LWX126}, proof of Lemma 7.16.)
\end{theorem}
\section{Permanence and exclusions used by the geometric classification}
\begin{theorem}[The known Kervaire squares and the high-square exclusion]
\label{frontier:Literature:552cd6818d8d}
\level{fine}
The following statements concern the specified standard classes in the fixed
$2$-primary sphere Adams spectral sequence.
\begin{enumerate}
\item (Adapted from \cite{MayThesis1964}, Appendix A and Theorem II.7.9.)
For every $0\leq j\leq3$, the square $h_j^2$ survives nontrivially to $E_\infty$.
\item (Adapted from \cite{MahowaldTangora}, Section 8; independently proved in
\cite{BMMS1986}, Chapter VI, Proposition 1.18(ii).)
The square $h_4^2$ survives nontrivially to $E_\infty$.
\item (Adapted from \cite{HHR}, Theorem 1.1, p.~2.)
For every $j\geq7$, the square $h_j^2$ does not survive nontrivially to $E_\infty$.
\end{enumerate}
\lean{KIP126.Literature.Standard.may_low_squares_permanent,KIP126.Literature.Standard.mahowald_tangora_h4_square_permanent,KIP126.Literature.Standard.hhr_high_squares_not_permanent}
% Blueprint note: These are the three square inputs also consumed directly by Geometry. Their fixed-model proofs remain registered Literature responsibilities. The detailed primary-source verification and its limitations are retained in the classical second-line theorem in this chapter. The positive geometric index is imposed by Browder, not by deleting the valid j=0 survival input.
\end{theorem}

\begin{theorem}[Mixed second-line products and their differentials]
\label{frontier:Literature:83726561f77b}
\level{fine}
Write $H_{i,j}=h_i h_j$ for products of the specified standard classes in the
fixed $2$-primary sphere Adams spectral sequence.
\begin{enumerate}
\item (Adapted from \cite{MayThesis1964}, Appendix A and Theorem II.7.9.)
The products $H_{0,2},H_{0,3},H_{2,4}$ survive nontrivially to $E_\infty$.
\item (Adapted from \cite{Mahowaldetaj}, Theorem 1', p.~249.)
For every $j\geq3$, the product $H_{1,j}$ survives nontrivially to $E_\infty$.
\item (Adapted from \cite{IWX}, Lemma 5.10(1), p.~143.)
There is a nonzero $d_3$ from $H_{2,5}$ at $(2,36)$ to degree $(5,38)$.
\item (Adapted from \cite{IWX}, Table 7, p.~225.)
There is a nonzero $d_4$ from $H_{3,5}$ at $(2,40)$ to degree $(6,43)$.
\item (Adapted from \cite{IWX}, Theorem 5.7, p.~142, and Table 6, p.~223.)
There is a nonzero $d_3$ from $H_{3,6}$ at $(2,72)$ to degree $(5,74)$.
\item (Adapted from \cite{IWX}, Table 6, p.~222.)
There is a nonzero $d_3$ from $H_{0,4}$ at $(2,17)$ to degree $(5,19)$.
\end{enumerate}
\lean{KIP126.Literature.Standard.may_two_line_products_permanent,KIP126.Literature.Standard.mahowald_eta_family_permanent,KIP126.Literature.Standard.mahowald_tangora_h2_h5_d3,KIP126.Literature.Standard.mahowald_tangora_h3_h5_d4,KIP126.Literature.Standard.iwx_h3_h6_d3,KIP126.Literature.Standard.iwx_h0_h4_d3}
% Blueprint note: These six inputs retain their existing fixed-model proof obligations. They form a separate compiled dependency group from the three square inputs used by Geometry. Source-to-classical comparison and primary-source qualifications are recorded with the classical second-line theorem in this chapter; no proof responsibility is removed by this regrouping.
\end{theorem}
